\documentclass[10pt]{article}

\usepackage[margin=0.75in, letterpaper]{geometry}
\usepackage{parskip}

\usepackage{amsthm, amsmath, amssymb}
\usepackage[pdftex]{graphicx}
\usepackage{hyperref}

\usepackage{enumerate} % For use of (a), (b), et cetera
\usepackage{booktabs} % Tables
\usepackage[margin=20pt, labelfont=bf, labelsep=period,
justification=justified]{caption} % Captions in figure floats

% https://stackoverflow.com/questions/3175105/inserting-code-in-this-latex-document-with-indentation
\usepackage{listings}
\usepackage{color}

\definecolor{dkgreen}{rgb}{0,0.6,0}
\definecolor{gray}{rgb}{0.5,0.5,0.5}
\definecolor{mauve}{rgb}{0.58,0,0.82}
\definecolor{lightgray}{rgb}{0.9,0.9,0.9}

% http://texdoc.net/texmf-dist/doc/latex/listings/listings.pdf
\lstset{
  backgroundcolor=\color{lightgray},
  frame=none,
%  language=python,
  aboveskip=3mm,
  belowskip=3mm,
  showstringspaces=false,
  columns=flexible,
  basicstyle={\small\ttfamily},
  numbers=none,
  numberstyle=\tiny\color{gray},
  keywordstyle=\color{blue},
  commentstyle=\color{dkgreen},
  stringstyle=\color{mauve},
  breaklines=true,
  breakatwhitespace=true,
  tabsize=4,
%  mathescape=false,
}


% ======================
% Document setup, layout
% ======================

% The following metadata will show up in the PDF properties
\hypersetup{
    colorlinks = true,
    urlcolor = magenta,  % Links to URLs
    linkcolor = blue,  % Links within PDF
    pdfauthor = {Torch collaboration},
    pdfkeywords = {FLASH},
    pdftitle = {Torch Manual - \today},
    pdfsubject = {},
    pdfpagemode = UseNone
}

% Don't indent paragraphs
\setlength\parindent{0em}

% Slightly more compact lines
\linespread{0.95}

% ===============
% Useful commands
% ===============
\newcommand*{\mt}{\mathrm}
\newcommand*{\unit}[1]{\;\mt{#1}}  % vemod.net/typesetting-units-in-latex
\newcommand*{\abt}{\mathord{\sim}} % tex.stackexchange.com/q/55701
%\newcommand*{\del}{\nabla}
\newcommand*\mean[1]{\bar{#1}}
\renewcommand{\vec}[1]{\boldsymbol{#1}}  % Bold vectors
\newcommand*{\tsup}{\textsuperscript}

% Sets, operators
\newcommand*{\ints}{\mathbb{Z}}
\newcommand*{\ptl}{\partial}
\newcommand*{\dl}{\nabla}
\newcommand*{\dtl}{\mathrm{d}}

\newcommand*{\Msun}{M_\odot}
\newcommand*{\sun}{\odot}
\newcommand*{\kB}{k_\mathrm{B}}

\begin{document}

% =======
% Titling
% =======
\begin{center}
A Torch Manual (ATM) \\
Last updated: \today \\
Started: April 25, 2019
\end{center}

\tableofcontents
\newpage

% ==================================
% Introduction
% ==================================
\section{Introduction}

\textbf{A Torch Manual} (ATM) (1) holds developer knowledge, (2) explains code
features in more depth than published literature, and (3) provides historical
context to show how \textbf{Torch} has evolved.  ATM gives references and
equations for code features, but should generally not explain physics.

ISM physics references:
\begin{itemize}
    \item Stahler \& Palla 2004 Star Formation
    \item Draine 2011 Physics of ISM/IGM
    \item Klessen \& Glover 2014 Physical Processes in the ISM,
        \href{https://arxiv.org/abs/1412.5182}{arXiv:1412.5182},
        \href{https://link.springer.com/chapter/10.1007\%2F978-3-662-47890-5\_2}{Springer PDF}.
\end{itemize}

\textbf{Warning:} ATM describes some features that are unmaintained and/or in
development.  Some features may not exist in the main Torch repository and may
not even work.

% =================================
% Version History of Torch
% =================================
\section{Version History of Torch}

\textbf{Torch} is distributed in the public repository:
\url{https://bitbucket.org/torch-sf/torch}

Old versions of the code exist in private repositories:
\begin{itemize}
    \item Aaron, Spring-Summer 2019: \url{https://bitbucket.org/aarontran/flash4.5-rad}
    \item Josh, Fall 2018 - Spring 2019: \url{https://bitbucket.org/joshuawall/flash4.5}
    \item Josh, Fall 2016 - Fall 2018: \url{https://bitbucket.org/joshuawall/flash4.2.2-rad}
    \item Juan+Christian, Spring 2013 - Spring 2016: \url{https://bitbucket.org/mordecai/flash4-ita-amnh-jcim}
    \item Christian, Spring 2013 - Spring 2014: \url{https://bitbucket.org/mordecai/flash4-ita-cb}
\end{itemize}

To get access to old repositories, contact the devs.

The version \texttt{flash4.2.2-rad} contains many add-ons and units accumulated
from the mid-2000s to present; many add-ons are NOT yet ported into the public
Torch repository.

% =================================
% What is added onto vanilla FLASH?
% =================================
\section{What is added to vanilla FLASH?}

Not a comprehensive list.

\begin{itemize}

\item
``Mtmod'' is called at restart by \texttt{IO\_readCheckpoint.F90}.
This is a custom edit to handle RNG state dump/read.
\textbf{TODO 2019 April 25 make a unit for this}

% In development - not pushed to any public repository yet.
%New unit 2019 July: AT adding \texttt{source/numericalTools/Random}
%to prevent constantly rewriting code for random number generation.
%Also allows user to choose a seed!

\item
Bouchut split solver (in flash4.2.2-rad).
Christoph Federrath has some comments on usage:
\url{http://flash.uchicago.edu/pipermail/flash-users/2014-May/001459.html}

\item
\texttt{Grid\_checkGridVar.F90} is located in
\texttt{flash4.2.2-rad/source/Grid}
or \texttt{flash4.5-rad/source/Grid\_checkVar/}.
This subroutine reports where fields are going wonky.
In flash4.2.2-rad, hooks for grid var debugging are included in:

\texttt{Driver/DriverMain/Unsplit/Driver\_evolveFlash.F90} \\
\texttt{Driver/DriverMain/Unsplit/Couple\_AMUSE/Driver\_evolveFlash.F90} \\
\texttt{Driver/DriverMain/Split/Couple\_AMUSE/Driver\_evolveFlash.F90}

In flash4.5-rad, hooks for grid var debugging are included in:

\texttt{Driver/DriverMain/Unsplit/Couple\_AMUSE/Driver\_evolveFlash.F90}

To use this functionality, specify \texttt{--defines=DEBUG\_GRIDVAR} in your
\texttt{setup} call.  Modify the Driver file to report on whatever grid
variable you desire.

\item
In \texttt{Driver\_computeDt} the guardcell filling is changed by josh due to
an error of zero dens on restarts...

\item
Modified Riemann solver behavior (Josh, Nov 2018):
``Added the ability for the code to mark regions where data
reconstruction falls back to 1st order as "shocked" regions
so that when using the HYBRID Reimann solver the solver will
also fall back to HLL instead of HLLD or HLLC.''
See change at:
\url{https://bitbucket.org/joshuawall/flash4.5/commits/bc708aa9d27dfc04741bb8f21c007b27cf8f033b#chg-source/physics/Hydro/HydroMain/unsplit/hy_uhd_dataReconstOneStep.F90}

\end{itemize}

% =============================
% More on build/compile process
% =============================
\section{More on build/compile process}

\subsection{AMUSE build snags}

AMUSE: we have run into problems with OpenMPI C++ bindings.
This is an issue for new versions of OpenMPI (4.x); you may have to build
OpenMPI from source with a configure flag to request bindings.  See:
\begin{itemize}
    \item \url{https://github.com/amusecode/amuse/pull/366}
    \item \url{https://github.com/joshuawall/amuse/commit/7ed7f2e45679587a43cb7b1f172297160f9bb820}
\end{itemize}

AMUSE in August 2019 (after HDF5 configure update) doesn't work quite right...
first you need to change the configure call for \texttt{with-hdf}
(at command line, do \texttt{./configure --help}).
Then, need to manually fix \texttt{HDF5\_LIBS} in \texttt{config.mk}.
%See Aaron's notes, 2019 August 13.

\subsection{FLASH build snags}

FLASH with OpenMPI v4.0.0 may throw errors like:
\begin{verbatim}
    undefined reference to `mpi_type_extent_'
    undefined reference to `mpi_address_'
    undefined reference to `mpi_type_struct_'
    undefined reference to `mpi_type_hvector_'
\end{verbatim}
Resolution was (early 2019) to use OpenMPI nightly build with
`--enable-mpi1-compatibility`.  It should be fixed in 4.0.1.
To address, see \url{https://github.com/open-mpi/ompi/issues/6114}

FLASH 4.5: when building with parallel HDF5 1.10.5, you may get the error:
\begin{verbatim}
    io_h5_report_xfer_method.c: In function ‘io_h5_report_xfer_method’:
    io_h5_report_xfer_method.c:105:26: error: ‘H5D_MPIO_FILTERS’ undeclared (first use in this function)
           if (global_cause & H5D_MPIO_FILTERS) {
\end{verbatim}
\texttt{H5D\_MPIO\_FILTERS} was removed in HDF5 1.10.2, apparently. See:\\
\url{https://github.com/hhoeflin/hdf5r/issues/97#issuecomment-422478916}

Resolution: directly modify:
\begin{verbatim}
    source/IO/IOMain/hdf5/io_h5_report_xfer_method.c
\end{verbatim}
Here is the diff:
\begin{verbatim}
    104a105
    > #if (H5_VERSION_LE(1,10,1))
    107a109,113
    > #else
    >       if (global_cause & H5D_MPIO_NO_COLLECTIVE_MAX_CAUSE) {
    >       ADD_CAUSE_TO_STR("H5D_MPIO_NO_COLLECTIVE_MAX_CAUSE ");
    >       }
    > #endif
\end{verbatim}
This fix is already applied to FLASH 4.6.0.

\subsection{Misc build notes}

HDF5: if you built HDF5 from source, rather than using "module load",
you can review your HDF5 configuration in the plaintext file:
\begin{verbatim}
hdf5-1.10.5/lib/libhdf5.settings
\end{verbatim}
which exists in your HDF5 build directory (i.e., the place that contains bin,
lib, include, etc.)

% =============================
% Parameters in FLASH and Torch
% =============================
\section{Parameters in FLASH and Torch}

\subsection{FLASH parameters}

\begin{itemize}
    \item If initializing sim w/ feedback star, choose a small \texttt{dtinit}
        because hydro will not take appropriate timestep for winds/radiation
        until the star begins barfing onto the grid.
\end{itemize}

\subsection{\texttt{bridge\_multiples.py} parameters}

\begin{itemize}
    \item Choose \texttt{cloud\_mass}, or equivalently your IMF high-mass
        cutoff, manually.
    \item Adjust \texttt{sink\_rad} to match your NXB setting,
        should be 16 for NXB=16.
\end{itemize}


% ===================
% Driver timestepping
% ===================
\section{Driver timestepping}

Torch modifies the default FLASH timestep calculation.  Here is a pseudo-code
summary of Josh's scheme to compute dt at the top, instead of bottom, of the
FLASH evolution loop.
\begin{lstlisting}[language=Fortran]
call Driver_init(...)

    call Driver_verifyInitDt(...)

        dr_dt = dtInit  ! only if not restart

call Driver_evolveFlash(...)

    call Driver_computeDt(...)

        dtNew = ...  ! Hydro, posdef, radtrans, all units; after MPI_AllReduce
        dr_dtAdvect = dtNew
        dtNew = min( dtNew, dtOld*dr_tstepChangeFactor )
        dtNew = min( max( dtNew, dr_dtMin ), dr_dtMax )

    dr_dt = dr_dtNew

    call dr_shortenLastDt(dr_dt, dr_simTime, dr_tmax, shortenedDt, 1)

    ! evolve units with dr_dt

    if (shortenedDt)
        dr_dt = dr_dtNew
    dr_dtOld = dr_dt
\end{lstlisting}

Old \texttt{bridge\_multiples.py} code used \texttt{dr\_dtAdvect} for hydro
timestep in brige loop, and used separate bridge loop controls (bridge had own
dtinit) to control timestep for e.g. simulations w/ feedback stars.

New refactored code uses \texttt{dr\_dt} as hydro timestep and does NOT include
separate bridge loop timestepping controls.

% =============
% Hydro lessons
% =============
\section{Hydro}

Remember to check your small numbers in flash.par (smallt, smallp, smallu,
etc.).  The flash.par provided with Torch has reasonable defaults, but the
FLASH defaults may not work as well.

Unsplit solver references:
\url{https://ui.adsabs.harvard.edu/abs/2009JCoPh.228..952L/abstract} \\
\url{https://ui.adsabs.harvard.edu/abs/2013JCoPh.243..269L/abstract}

MHD code comparison:
\url{https://ui.adsabs.harvard.edu/abs/2011ApJ...737...13K/abstract}

\subsection{CFL and pressure holes}

A problem that has recurred through work with Juan, Josh, etc.
Sometimes, lowering CFL to push past can help.
Bouchut, unsplit, \texttt{use\_PosdefComputeDt} are supposed to help alleviate.
Helpful thread:
\url{http://flash.uchicago.edu/pipermail/flash-users/2019-May/002908.html}

% --------------------
% Artificial viscosity
% --------------------
\subsection{Artificial viscosity}

About 0.1 is good.  See Fryxell paper section 3.2 for brief description, I am
unsure if the Lapidus paper actually matches whatever is in the code.

Only turns on when shocks are present.  For example, using the 1-D advection
equation for just density:
\[
    \frac{\ptl \rho}{\ptl t} + \frac{\ptl}{\ptl x} \left( \rho u \right) = 0
\]
where $\rho u$ is the flux, the artificial viscosity looks something like
\[
    \frac{\ptl \rho}{\ptl t}
    + \frac{\ptl}{\ptl x} \left(
        \rho u - C \frac{\ptl u}{\ptl x} \frac{\ptl \rho}{\ptl x} \left(\Delta x\right)^2
    \right)
    = 0
\]
where $C$ is the dimensionless runtime parameter \texttt{avisc},
and the viscosity is only turned on in cells satisfying $\ptl u / \ptl x < 0$;
i.e., zones that are steepening.
This all based on reading \texttt{unsplit/hy\_uhd\_getFaceFlux.F90}

It looks consistent with artificially smearing the flux, EXCEPT that you
only keep the cross-term for some reason . . . why this exact form?
\[
    \frac{\ptl^2}{\ptl x^2} \left( \rho u \right)
    = \rho \frac{\ptl^2 u}{\ptl x^2}
    + u \frac{\ptl^2 \rho}{\ptl x^2}
    + 2 \frac{\ptl \rho}{\ptl x} \frac{\ptl u}{\ptl x}
\]

% =============
% Gravity
% =============
\section{Gravity}

% ------------------------------
% Gravity - unsplit hydro solver
% ------------------------------
\subsection{Unsplit Gravity Calculation}

The unsplit MHD solver computes gravity in \texttt{hy\_uhd\_unsplit.F90}.
Hydro is evolved first in each timestep.  Then, we must update grav potential
and accel for other evolution subroutines (e.g., particle advance).

\begin{lstlisting}[language=Fortran]
! compute GPOT_VAR from Poisson equation via Gravity_potentialListOfBlocks(...)
! also compute sink->gas accel SGA{X,Y,Z}_VAR and gas->sink accel ACC{X,Y,Z}_PART_PROP
call hy_uhd_prepareNewGravityAccel(blockCount,blockList,...)

do i = 1,blockCount
    ! compute grav{X,Y,Z} accels via Gravity_accelOneRow(...),
    ! accel is SGA{X,Y,Z}_VAR plus finite difference of GPOT_VAR
    call hy_uhd_putGravityUnsplit(blockID,...,gravX,gravY,gravZ)

    ! update velocities and energy from grav{X,Y,Z}
    call hy_uhd_addGravityUnsplit(blockID,...,gravX,gravY,gravZ,...)
end do
\end{lstlisting}

\textbf{Note:} in normal FLASH, \texttt{GPOT\_VAR} includes contribution
from massive non-sink particles via \texttt{PDEN\_VAR}.  This is disabled in
\textbf{Torch} since all massive non-sink particles are handled by the AMUSE
bridge.

\textbf{Question:} are sink particles double counted in grav accel?  They require the
\texttt{active/massive/Leapfrog} integrator, which defines both
\texttt{MASS\_PART\_PROP} and \texttt{PDEN\_VAR}.  Superficially, it may appear
that sinks are included twice:
\begin{itemize}
\item \texttt{Gravity\_potentialListOfBlocks}: \texttt{GPOT} includes sink mass via \texttt{PDEN}, and
\item \texttt{Particles\_sinkAccelGasOnSinksAndSinksOnGas}: computes sink acceleration upon gas directly.
\end{itemize}

\textbf{Answer:} no, we are OK.
\begin{itemize}
\item In default FLASH: in \texttt{Particles\_updateGridVar}, which would map
prtl masses to PDEN, a special exclusion is made for sink particles.  Only
non-sink massive prtls are included in PDEN.
\item In Torch: all massive prtls excluded from GPOT.
\end{itemize}

% --------------------------
% Gravity - sink particles
% --------------------------
\subsection{Sink particle gravity}

In \texttt{Particles\_sinkAdvanceParticles} sink particles are advanced using
both gas $\to$ sink and sink $\leftrightarrow$ sink accelerations.
The gas $\to$ sink acceleration \texttt{ACC{X,Y,Z}\_PART\_PROP} is not re-computed,
but just taken and assumed up-to-date, which I think presumes?? that the Hydro
unit calculated it for us??

The sink $\leftrightarrow$ sink acceleration is computed by direct calculation,
according to FLASH manual.

% --------------------------
% Gravity - static potential
% --------------------------
\subsection{Static Grav Potential}

Used in stratified box simulations to get disk grav, exponential z-profile
(Joung, Bacyznski, Ibanez-Mejia work).

Use \texttt{flash4.2.2-rad} with BHTree, but not BHTree4.3.
The FLASH source code is directly modified in:
\begin{verbatim}
    flash4.2.2-rad/source/physics/Gravity/GravityMain/Poisson/Gravity_accelOneRow
\end{verbatim}
to hook in a static grav acceleration.
But, BHTree4.3 will override this.
All AMUSE couplings in \texttt{flash4.2.2-rad} will also override.

%In addition, sink particles do NOT feel this grav potential!

% ==============
% Gravity Bridge
% ==============
\section{Gravity Bridge}

Pseudocode for stars to kick gas+sinks, AMUSE bridge.

\begin{lstlisting}[language=Fortran]
FUNCTION get_gravity_particles_on_gas(dt, kick_number)

    if (kick_number == 1) then

        VEL{X,Y,Z}_VAR += BGA{X,Y,Z}_VAR * dt  ! gas kicked by stars
        BGA{X,Y,Z}_VAR = 0

    else if (kick_number == 2) then

        call Gravity_potentialListOfBlocks(GPOT_VAR)
        ! Result: GPOT_VAR computed from DENS_VAR (solve Poisson eqn).
        !          sinks are NOT included in GPOT_VAR, a change from default FLASH
        ! Side effect: update SGA{X,Y,Z}_VAR = grav accel field due to sinks
        ! Should this be in get_gravity_gas_on_particles instead? -AT,2020feb27

        call Gravity_potentialListOfBlocksPDENonly(BGPT_VAR, ACTIVE_PART_TYPE)
        ! result: BGPT_VAR computed from PDEN_VAR (solve Poisson eqn)
        !          where PDEN_VAR = prtl-to-mesh mapping of stars

        call Gravity_accelListOfBlocks(IAXIS, BGAX_VAR, BGPT_VAR)
        call Gravity_accelListOfBlocks(JAXIS, BGAY_VAR, BGPT_VAR)
        call Gravity_accelListOfBlocks(KAXIS, BGAZ_VAR, BGPT_VAR)
        ! result: BGA{X,Y,Z}_VAR contains {x,y,z}-acceleration from BGPT_VAR

        VEL{X,Y,Z}_VAR += BGA{X,Y,Z}_VAR * dt  ! gas kicked by stars

    end if

    call Particles_longRangeForce(SINK_PART_TYPE, BGPT_VAR)
    ! result: sink ACC{X,Y,Z}_PART_PROP contains acceleration from BGPT_VAR

    VEL{X,Y,Z}_PART_PROP += ACC{X,Y,Z}_PART_PROP * dt  ! sinks kicked by stars

END FUNCTION
\end{lstlisting}

Pseudocode for gas+sinks to kick stars, AMUSE bridge.

\begin{lstlisting}[language=Fortran]
FUNCTION get_gravity_gas_on_particles(dt)

    ! pre-requisite: SGA{X,Y,Z}_VAR = sink accel field is up-to-date
    Particles_longRangeForce(ACTIVE_TYPE_PART, GPOT_VAR)
    ! result: star ACC{X,Y,Z}_PART_PROP holds accel from gas (GPOT_VAR) and sinks (SGA{X,Y,Z}_VAR)

    ! breakdown of Particles_longRangeForce(...)
    !   call Gravity_accelListOfBlocks(IAXIS, GRAC_VAR, GPOT_VAR)
    !       for all blocks:
    !           call Gravity_accelOneRow(SWEEP_X, GPOT_VAR)
    !               returns: first-order finite difference of GPOT_VAR along x, plus SGAX_VAR
    !       result: GRAC_VAR contains x-acceleration from gas (GPOT_VAR) and sinks (SGA{X,Y,Z}_VAR)
    !   call Grid_mapMeshToParticles(...)
    !       result: GRAC_VAR mapped to ACCX_PART_PROP via CIC method
    !   repeat for other coordinates JAXIS, KAXIS

    VEL{X,Y,Z}_PART_PROP += ACC{X,Y,Z}_PART_PROP * dt  ! stars kicked by gas+sinks

END FUNCTION get_gravity_gas_on_particles
\end{lstlisting}

NOTE: the coupling code subroutines:
\begin{verbatim}
    Gravity_getAccelAtPoint.F90
    Gravity_getPotentialAtPoint.F90
\end{verbatim}
are not currently in use.  Josh used these in an older version of bridge
interface, but all that stuff is commented out now.  These routines are not
called anywhere in FLASH code.
- AT, 2019 May 09.


% ==========================
% Particles: sinks and stars
% ==========================
\section{Particles: sinks and stars}

% --------------------------------------------
% Particle data structures, generally speaking
% --------------------------------------------
\subsection{Particle data structures}

\texttt{Particles\_data.F90} provides:
\begin{verbatim}
    particles(:,:)      (NPART_PROPS, MAX_PARTICLES_PER_PROCESSOR=pt_maxPerProc)
    pt_typeInfo(:,:)    (PART_TYPE_INFO_SIZE, NPART_TYPES)
\end{verbatim}

All properties are shared by all particle types.  The \texttt{particles} is
manually sorted by prtl type, then by blk, at various points in the FLASH
evolution loop.

\texttt{pt\_typeInfo} holds metadata about each particle type.
These metadata properties are declared in \texttt{Particles.h}, they may include some of:
\begin{verbatim}
    PART_TYPE = type of particle
    PART_MAPMETHOD = method to map to/from mesh
    PART_TYPE_BEGIN = starting index of particle type in particles(:,:) array
    PART_INITMETHOD
    PART_LOCAL = local number of this type of particle
    PART_LOCAL_MAX = maximum of this type allowed on proc
    PART_ADVMETHOD
    PART_INTEGRATE_METHOD = time integration algorithm
\end{verbatim}
Example: to iterate over the stars (active particles), you can do:
\begin{lstlisting}[language=Fortran]
    p_begin = pt_typeInfo(PART_TYPE_BEGIN,ACTIVE_PART_TYPE)
    p_num = pt_typeInfo(PART_LOCAL,ACTIVE_PART_TYPE)  ! number of stars on current proc
    do p = p_begin, p_begin+p_num-1
        print*, "Particle mass =", particles(MASS_PART_PROP, p)
    end do
\end{lstlisting}

\texttt{Particles\_sinkData.F90} provides:
\begin{verbatim}
    particles_local(:,:)        (pt_sinkParticleProps, sink_maxSinks)
    particles_global(:,:)       (pt_sinkParticleProps, sink_maxSinks)
\end{verbatim}


% ------------------------
% Star/sink data structures
% ------------------------
\subsection{Star/sink data structures}

The AMUSE interface to FLASH uses pointers to get/set different particle types.

The sink and star data reside in:
\begin{verbatim}
Particles/ParticlesMain/active/Sink/Couple_AMUSE/Particles_sinkData.F90
Particles/ParticlesMain/active/Sink/Couple_AMUSE/Couple_AMUSE_Sinks_and_Stars/Particles_data.F90
\end{verbatim}
respectively.

AMUSE interface: \texttt{set\_particle\_pointers} tells the AMUSE worker, OK,
here are the particles you can access directly from memory!

When AMUSE interface code points to star particles:
\begin{verbatim}
    AMUSE interface.F90                 FLASH Particles_data.F90
    -------------------                 ------------------------
    real :: particles_pointer(:,:)  =>  particles(:,:)
    int :: num_part_local_prt       =>  pt_numLocal
    int :: new_particles_tags(:)    =>  new_massive_tags(:)
    int :: number_new_particles     =>  number_new_massive
\end{verbatim}

When AMUSE interface code points to sink particles:
\begin{verbatim}
    AMUSE interface.F90                 FLASH Particles_sinkData.F90
    -------------------                 ------------------------
    real :: particles_pointer(:,:)  =>  particles_local(:,:)
    int :: num_part_local_prt       =>  localnp
    int :: new_particles_tags(:)    =>  new_sink_tags(:)
    int :: number_new_particles     =>  number_new_sinks
\end{verbatim}

Note \texttt{particles\_pointer} has array shape
\texttt{(NPART\_PROPS, MAX\_PARTICLES\_PER\_PROCESSOR)}.

%AMUSE interface: \texttt{get\_particle\_tags} and \texttt{get\_new\_tags} do
%MPI gather to get all particle tags.


% -----------------------------------
% Sink creation, accretion, evolution
% -----------------------------------
\subsection{Sink creation, accretion, evolution}

Mostly summarized by Federrath
\href{https://ui.adsabs.harvard.edu/abs/2010ApJ...713..269F/abstract}{2010ApJ...713..269F}

Sink motion is a two-step process.
First, we compute gas-to-sink and sink-to-sink gravity, updating the sink
positions and velocities in a leapfrog-staggered manner.
Second, we check whether sink particles need to be created, and then accrete
gas onto all existing and newly created sink particles.

The first step happens in \texttt{Particles\_sinkAdvanceParticles}, which is
called by \texttt{Particles\_advance} during the \texttt{Driver\_evolveFlash}
loop.  Below is stripped-down pseudocode to show process for the $i$-th sink
particle along the $x$ direction only.
\begin{lstlisting}[language=Fortran]
! store acceleration from GasOnSinksAccel (now called in Gravity_PotentialListOfBlocks)
ax_gas(i) = particles_local(ACCX_PART_PROP, i)

do while (.not. end_subcycling)

   particles_local(ACCX_PART_PROP, i) = ax_gas(i)

   ! compute sink-sink gravity
   call pt_sinkAccelSinksOnSinks(local_min_radius, local_max_accel)

   ! First step of leapfrog:
   particles_local(VELX_PART_PROP, i) += 0.5*dt * particles_local(ACCX_PART_PROP, i)
   particles_local(POSX_PART_PROP, i) += dt * particles_local(POSX_PART_PROP, i)

   ! restore gas contribution
   particles_local(ACCX_PART_PROP, i) = ax_gas(i)

   ! re-compute sink-sink acceleration at new prtl location
   call pt_sinkAccelSinksOnSinks(local_min_radius, local_max_accel)

   ! Second step of leapfrog:
   particles_local(VELX_PART_PROP, i) += 0.5*dt * particles_local(ACCX_PART_PROP, i)

end do
\end{lstlisting}

The next step happens in \texttt{Particles\_sinkCreateAccrete}, with the
following logic.

\begin{lstlisting}[language=Fortran]
! loop over all LEAF BLOCKS (Federrath Sec 2.2)
do lb = 1, blockCount
   do kp, jp, ip = blklimits(...)
      if (comovingCellDens .gt. density_thresh_comoving) then
         ! 1. check for max refinement  (JOSH added/fixed??? this???)
         ! 2. check for neighbor sinks
         ! 3. check converging flow in adjacent (3x3x3) cells
         ! 4. check grav potential minimum
         ! 5. check jeans mass
         ! 6. check total energy < 0, exclude bulk COM KE.
   end do
end do

! loop over all LEAF BLOCKS with dens > threshold (Federrath Sec 2.3)
do llb = 1, blockCount
   do kp, jp, ip = blklimits(...)

      mass = (comovingCellDens - density_thresh_comoving) * dVol

      ! does current cell lie within accretion radius of any sinks?
      if (np_found .gt. 0) then

         do npf = 1, np_found
            radius(npf) = ... ! dist from cell to sink
            etot(npf) = ...  ! grav + kinetic energy, computed including only the mass above density thresh
            vr(npf) = ...  ! radial velocity
         end do

         inner_r_accr = max( 0.2*accretion_radius_comoving, (dVol**(1.0/3.0)) )

         do npf = 1, np_found

             ! FLASH 4.5 VERSION, WITHOUT ANDREA GATTO BUG FIX
             !   ...
             !   IF cell lies within "inner" accretion radius
             !   of any sink(s), always accrete to most
             !   strongly bound sink
             !   ...
             !   ELSE, check Etot<0, vr < 1e-5*cs,
             !   if satisfied, then accrete to most strongly
             !   bound sink

             ! ANDREA GATTO BUG FIX
             !   ...
             !   IF cell lies within "inner" accretion radius
             !   of any sink(s), query Etot
             !   ...
             !   ELSE, check Etot<0, vr < 1e-5*cs,
             !   if satisfied, query Etot
             !   ...
             !   RESULT, take lowest Etot cell that is
             !   either "inner" or ("outer" + bound)
             !   this allows cell to accrete to a sink
             !   outside "inner" radius if it is more
             !   strongly bound to that sink AND
             !   passes the Etot, vr checks

         end do  ! loop over all nearby sinks
      end if  ! cell is near >=1 sink
   end do  ! coords
end do  ! blocks

call pt_sinkGatherGlobal(...) ! MPI communicate sink properties

do lp = 1, localnp
   ! check if local particle is affected by regions on other CPUs
   do nlp = localnp+1, localnpf
      if(int(particles_local(iptag,lp)) .eq.  int(particles_global(iptag,nlp))) then
         tot_mass(lp) = tot_mass(lp) + particles_global(ipm,nlp)
         ...
      end if
   end do

   ! update particle properties (conservation laws)
   particles_local(iold_pmass,lp) = particles_local(ipm,lp)
   if (tot_mass(lp) .ne. 0.0) then
       ! mass update
       particles_local(ipm,lp) = particles_local(ipm,lp) + tot_mass(lp)
       ! position update
       px_old = particles_local(ipx,lp)
       py_old = particles_local(ipy,lp)
       pz_old = particles_local(ipz,lp)
       particles_local(ipx,lp) = particles_local(ipx,lp) + cm_x(lp)/particles_local(ipm,lp)
       particles_local(ipy,lp) = particles_local(ipy,lp) + cm_y(lp)/particles_local(ipm,lp)
       particles_local(ipz,lp) = particles_local(ipz,lp) + cm_z(lp)/particles_local(ipm,lp)
       ...
   end if
end do

! Josh added queries to save local gas conditions as sink properties
\end{lstlisting}

After accreting, the sink immediately adopts the center-of-mass (COM) position
and velocity of the combined sink-and-accreted-gas system.

The COM position/velocity update includes the sink's own mass in a subtle way.
The sink's own mass enters via total mass \texttt{particles\_local(ipm,lp)}
during the position update step.

Andrea Gatto bug fix... motivation unknown to AT (2019 dec 18).
Josh picked it up from Andrea and added it at some point.

% -----------
% Stellar IMF
% -----------
\subsection{Stellar IMF sampling}

%Key questions: is there a high-mass cutoff?  What is low-mass slope?  Does IMF
%for clusters differ from galaxy-integrated IMF?

The main star-forming algorithm uses a Kroupa IMF, with a marginally different
index than the Kroupa papers cited:
\begin{align*}
    \frac{\mathrm{d} N}{\mathrm{d} M} \propto
    \begin{cases}
        M^{-0.30}, & 0.01 < M/\Msun < 0.08 \\
        M^{-1.30}, & 0.08 < M/\Msun < 0.50 \\
        M^{-2.30}, & 0.50 < M/\Msun
    \end{cases}
\end{align*}


Salpeter (\href{https://ui.adsabs.harvard.edu/abs/1955ApJ...121..161S}{1955ApJ...121..161S}) IMF:
\[
    \frac{\mathrm{d} N}{\mathrm{d} M} \propto M^{-2.35}
\]
for $-0.4 < \log_{10} \left(M/\Msun\right) < +1.0$, or roughly $0.4 \Msun$ to
$10 \Msun$.

Weidner and Kroupa
\href{http://adsabs.harvard.edu/abs/2004MNRAS.348..187W}{2004MNRAS.348..187W}
Section 2 states:
\begin{align*}
    \frac{\mathrm{d} N}{\mathrm{d} M} \propto
    \begin{cases}
        M^{-0.30}, & 0.01 < M/\Msun < 0.08 \\
        M^{-1.30}, & 0.08 < M/\Msun < 0.50 \\
        M^{-2.30}, & 0.50 < M/\Msun < 1.00 \\
        M^{-2.35}, & 1.00 < M/\Msun
    \end{cases}
\end{align*}

Weidner, Kroupa, Pflamm-Altenburg
\href{http://adsabs.harvard.edu/abs/2013MNRAS.434...84W}{2013MNRAS.434...84W}
Appendix B states:
\begin{align*}
    \frac{\mathrm{d} N}{\mathrm{d} M} \propto
    \begin{cases}
        M^{-0.30}, & 0.01 < M/\Msun < 0.08 \\
        M^{-1.30}, & 0.08 < M/\Msun < 0.50 \\
        M^{-2.35}, & 0.50 < M/\Msun < M_\mathrm{max}
    \end{cases}
\end{align*}


% ----------------
% Multiples
% ----------------
\subsection{Multiples}

In AMUSE's multiples module, a bound binary is represented by one
center-of-mass (COM) particle and two star particles (leaves).  For now,
consider only binaries and ignore triples or higher-order systems.

The COM particle is the ``root'' of its tree.  It exists in FOUR different data
structures:
\begin{verbatim}
    grav.particles              # ph4 buffer
    mult._inmemory_particles    # copy (not view) of ph4 buffer
    mult.root_to_tree = {
        root: tree,             # the key "root" is a particle
        ...
    }
    tree.particle               # taken from root_to_tree
\end{verbatim}
The four particle copies are usually updated together.

The \texttt{tree} is an instance of:
\begin{verbatim}
    amuse.datamodel.trees.BinaryTreeOnParticle
\end{verbatim}
which represents (1) a node and two children.  The children may in turn be COM
particles, and so on and so forth.

Warning: in Multiples, the \texttt{root} and \texttt{tree.particle} are usually
NOT synced with gravity code.  The \texttt{tree.particle} captures the system
state when the binary first formed.  However, to account for stellar evolution,
\texttt{TORCH} modifies and update these particles (and the leaf particles /
stars as well).

Warning: the AMUSE logic seems to assume \textbf{full} binary trees; that is,
every node has either zero or two children.  Some of its functions will NOT
work correctly for trees with one-child nodes.

Historical note: Josh experimented with channels to Multiples.  It seems like
this did not work, since the \texttt{mult.stars} Particles set is constructed
dynamically from the \texttt{mult.root\_to\_tree} and probably doesn't play
well with AMUSE channels.  Josh settled for iterating through the Multiples
tree at every sync operation to Multiples.

%AT has hand-written notes from study/reading of Multiples, from late 2019 October.


% ==============================================
% Ray tracing, ionization, heat/cool sub-cycling
% ==============================================
\section{The Sub-Cycle: rays, ionization, heat/cool}

\textbf{One moral:} in runtime parameters, set the ionization fraction $x$ to the
equilibrium value for your initial temperature and density.
Cold gas can take $> 10^9 \unit{yr}$ for cosmic ray ionization and
recombination to balance; if the initial ionization is wrong, the heating and
cooling may be wrong too.

Timestep sub-cycling: only triggered for rapidly varying radiation fields.
Tentatively (as of April 2019), heating and cooling is NOT sub-cycled.
IF the ionization is stable, \texttt{RadTrans} will take two steps.
It first computes an ionization timestep using \texttt{rt\_dt = 0.01 dt}.
Then, it advances rays, ionization, and heating/cooling using (almost) the full
hydro timestep, \texttt{rt\_dt = 0.99 dt}.

Ionization: implicit solve + timestep-estimation is explained in Josh's thesis
and Paper II.  Baczynski used an explicit solve and accounted for electrons
from metals.

\textbf{TODO low priority (2019 April 25):} when the gas is almost fully
ionized and $1-x \ll 1$, the quadratic solve returns $x$, but the value of
$1-x$ may be inaccurate (subtracting two big numbers to get a small number).
Not tested, not yet demonstrated whether this is a problem or not.

\textbf{TODO low priority (2019 May 01):} do the ionization more simply when
there are no stars.  No need for clever sub-cycling with rad transfer in this
case, just advance the ionization and proceed!

% ===========
% Ray tracing
% ===========
\section{Ray tracing}

The code is split across two units.  First a particle unit:
\texttt{Particles/ParticlesMain/ray\_pe}
then a rad transfer unit:
\texttt{physics/RadTrans/RadTransMain/RayRad} .
Two PPDEFINE flags indicate presence of ray tracing units:
FERVENT and PE\_HEATING.  Both ray and ray\_pe versions of unit will define
FERVENT.

There exists an old version of the raytracing particle unit,
\texttt{Particles/ParticlesMain/ray}, that HAS been modified by Josh; not sure
if it's up to date or if it works.


% ==========
% Ionization
% ==========
\section{Ionization}

\textbf{Torch} solves an ionization equation for a pure hydrogen gas:
\begin{equation}
    \frac{\dtl x}{\dtl t}
    = C_\mt{cl}(T) n_\mt{H} x (1-x) + k_\mt{ion} (1-x) - \alpha_B(T) n_\mt{H} x^2
\end{equation}
or equivalently
\[
    \frac{\dtl n_e}{\dtl t}
    = C_\mt{cl}(T) n_e n_\mt{HI} + k_\mt{ion} n_\mt{HI} - \alpha_B{T} n_e n_p
\]
in \texttt{calc\_ionization.F90}.
The ionization fraction $x = n_e/n_\mt{H}$, the total H nucleon density
$n_\mt{H} = n_\mt{HI} + n_e$, and the electron/proton densities $n_e$ and $n_p$
are equal.  Compared to FERVENT, \textbf{Torch} omits electrons from metals,
dust-catalyzed recombination, and CR ionization of molecules.

\subsection{Ionization Terms}

Collisional ionization follows equation (3) of Cox and Tucker
\href{1969ApJ...157.1157C}{http://adsabs.harvard.edu/abs/1969ApJ...157.1157C}.
Draine's textbook restates and abridges Cox/Tucker's result.
\begin{align}
    C_\mt{cl}
    &= \frac{1.3 \times 10^{-8} \times 0.83 \times 1 \unit{cm^{3}\;s^{-1}}}{13.598^2}
    \;
    T^{1/2} \exp \left( - \frac{13.598 \unit{eV}}{\kB T}\right) \\
    %&= \left( 5.8 \times 10^{-11} \unit{cm^{3}\;s^{-1}} \right)
    %T^{1/2} \exp \left( - \frac{13.598 \unit{eV}}{\kB T}\right)
    &= 5.8 \times 10^{-9} \unit{cm^{3}\;s^{-1}}
    \left(\frac{T}{10^4\unit{K}}\right)^{1/2} \exp \left( - \frac{13.598 \unit{eV}}{\kB T}\right)
\end{align}

The ``Case B'' recombination coefficient is from Draine's book; it is good for
$T = 30$ to $30,000\unit{K}$:
\[
    \alpha_B
    = 2.54 \times 10^{-13} \unit{cm^{3}\;s^{-1}}
    \left( \frac{T}{10^4 \unit{K}} \right)^{-0.8163 - 0.0208\ln(T/10^4\unit{K})}
\]
C.~Baczynski uses a simpler expression that is credited to Osterbrock (1989).
The Osterbrock and Draine expressions differ by $\abt 20\%$ or less over the
temperature range given by Draine.
\[
    \alpha_B
    = 2.59 \times 10^{-13} \unit{cm^{3}\;s^{-1}}
    \left( \frac{T}{10^4 \unit{K}} \right)^{-0.7}
\]

\begin{figure}
    \center
    \includegraphics[width=0.5\textwidth]{{fig/recombination_case_b}.pdf}
    \caption{
        Recombination case B coefficients.
    }
    %\label{fig:peheat}
\end{figure}

The photoionization rate $k_\mt{ion}$ combines (1) background cosmic-ray (CR)
ionization of \textit{molecular} H2 to H2+, and (2) ionizing photons from FUV-emitting
stars.  The CR ionization is from eqn (27) and Table 4 of Padovani, Galli, and
Glassgold
\href{https://ui.adsabs.harvard.edu/abs/2009A\%26A...501..619P/abstract}{2009A\&A...501..619P}.
Josh uses ``W98'' spectrum shown in Figure 13.
\textbf{WARNING/TODO:} the exponential cutoff for $N(H_2) > 10^{25}
\unit{cm^{-2}}$ is NOT IMPLEMENTED.  For that, see the paper's corrigendum:
\href{https://ui.adsabs.harvard.edu/abs/2013A\%26A...549C...3P/abstract}{2013A\&A...549C...3P}.
\[
    k_\mt{ion,CR} = \zeta^{H_2}_k
    = 2 \times 10^{-17} \unit{s^{-1}} \left( \frac{N(H_2)}{10^{20} \unit{cm^{-2}}} \right)^{-0.021}
\]
To convert from volume density to column density, Josh uses a local
approximation valid for $N(H_2) \sim 10^{20}$ to $10^{25} \unit{cm^{-2}}$:
\begin{align*}
    N(H_2) &\approx \lambda_\mt{J} n_\mt{H} \\
    \lambda_\mt{J} &= \sqrt{\frac{5 \pi \kB T}{\mu \rho G}} \propto c_s t_\mt{ff}
\end{align*}
where $\mu = \rho / n$ is mean particle mass.

%\textbf{TODO 2019 April 25:} column density approximation, do we need to halve
%the number density from FLASH?  Relates to assumptions on molecular vs atomic
%gas fraction.  Very unimportant here: the CR ionization rate is almost constant
%with respect to number density.  But may matter elsewhere.

\textbf{TODO 2019 May 31:} Walch+ 2015 (SILCC) mention: ``at no point do we
assume that hydrogen is in collisional ionization equilibrium: its contribution
to the cooling rate is determined self-consistently at all temperatures,
allowing us to properly model the non-equilibrium effects that can strongly
affect the cooling rate at temperatures close to $10^4$ K''.
\textbf{Do we need to consider this?  Need to read more}

\subsection{Ionization Equilibrium}

To find the correct starting ionization for your simulation, set $\dtl x/\dtl t
= 0$.  This gives a quadratic for the equilibrium value of $x$:
\begin{equation}
    0 = \left( C_\mt{cl}(T) + \alpha_B(T) \right) n_\mt{H} x^2
    + \left(-C_\mt{cl}(T) + k_\mt{ion}\right) x
    - k_\mt{ion}
\end{equation}
Beware, the solution is not so accurate for $x \approx 1$.
I find unphysical $x > 1$ when solving this quadratic with $\rho = 10^{-25}
\unit{g\;cm^{-3}}$ and $T = 10^6 \unit{K}$.
It is probably necessary to swap variable to $1-x$ to do the calculation
accurately.


% ===============
% Heating/Cooling
% ===============
\section{Heating/Cooling}

Like ionization, heat/cool is intimately tied to radiation transfer.
Location: \\
\texttt{HeatMain/HeatCool/phenHeat/mol\_and\_dust/solver}.

\textbf{Developer warning:} lots of global variables used here.

Truth table for heating/rad parameters:
\begin{verbatim}
    useHeat  rt_useRadTrans  rt_heatInRad
    F        ...             ...            Heat.F90 does nothing
    T        F               F              Heat.F90 calls RadHeat
    T        F               T              Heat.F90 calls RadHeat
    T        T               F              Heat.F90 calls RadHeat
    T        T               T              Heat.F90 does nothing, RadTrans calls RadHeat
\end{verbatim}

There is also the parameter \texttt{useRadTransfer}, which toggles the code for
ray tracing in \texttt{ParticlesMain/ray} and \texttt{ParticlesMain/raype};
therefore \texttt{useRadTransfer} does not
directly affect heating/cooling in the above truth table! (but of course if you
turn off stars' radiation, that will change feedback heating)


% -------------------
% Gas Heating/Cooling
% -------------------
\subsection{Gas heat/cool}
obeys the equation
\begin{equation}
    \frac{\dtl \epsilon}{\dtl t} =
    \Gamma_\mt{phen} + \Gamma_\mt{CR} + \Gamma_\mt{PE}
    - \Lambda_\mt{atom} - \Lambda_\mt{mol} - \Lambda_\mt{g\to d}
\end{equation}
where $\epsilon$ is specific energy in units $\mathrm{erg\;g^{-1}}$,
so all $\Gamma$ and $\Lambda$ have units $\mathrm{erg\;s^{-1}\;g^{-1}}$.

\textbf{Method:} use implicit finite-difference scheme.
Let $f\left(\epsilon\right) = \dtl\epsilon/\dtl t$.  Then, the scheme is:
\[
    \frac{\epsilon^{n+1} - \epsilon^{n}}{\Delta t}
    =
    f \left(\epsilon^{n+1}, n^n, T_\mt{dust}^{n+1}, \ldots\right)
\]
where $n$ and $n+1$ index hydro timesteps or a heat/cool subcycle of size
$\Delta t$.  The code seeks the root of:
\[
    \epsilon^{n+1} - \Delta t f \left(\epsilon^{n+1}, \ldots\right) - \epsilon^n = 0
\]
\textbf{Developer warning:} the bracket- and root-finding code is called as follows:
\begin{lstlisting}[language=Fortran]
    call zbrac(f, x1, x2, xold, dt, success)
    call zero(a, b, a_old, machep, t, f, dt)
\end{lstlisting}
These functions do NOT operate on function \texttt{f}, but rather operate on the
left-hand side of the equation:
\begin{lstlisting}[language=Fortran]
    x - dt * f(dt, x) - xold = 0
\end{lstlisting}
in order to implement the implicit solve.  This is very confusing, but works
correctly.
Fortran functions are not first-class
functions; we cannot easily pass them around like objects.  So engineering the
bracket- and root- finding code in a general way is a bit non-trivial.  Maybe
modern Fortran features can alleviate?

\textbf{Numerics warning:} implicit solve is formally stable, but can still be
inaccurate.  See Townsend \url{http://adsabs.harvard.edu/abs/2009ApJS..181..391T}
for a discussion and an alternative scheme (not implemented in Torch, yet).

% ---------------------------
\subsubsection{Photo-heating}

$\Gamma_\mt{phen}$ is ``direct'' radiation heating (also called
photo-heating), \texttt{PHHE\_VAR/DENS\_VAR}.

% ------------------------
\subsubsection{CR heating}

$\Gamma_\mt{CR}$ is heating from cosmic rays, eqn (27) and Table 4 of
Padovani, Galli, and Glassgold
\href{https://ui.adsabs.harvard.edu/abs/2009A\%26A...501..619P/abstract}{2009A\&A...501..619P}.
\[
    \Gamma_\mt{CR} =
        \left(20 \unit{eV}\right) \left(2 \times 10^{-17} \unit{s^{-1}}\right)
        \left( \frac{N(H_2)}{10^{20} \unit{cm^{-2}}} \right)^{-0.021}
    =   6.41 \times 10^{-28} \unit{erg\;s^{-1}}
        \left( \frac{N(H_2)}{10^{20} \unit{cm^{-2}}} \right)^{-0.021}
\]
where $20 \unit{eV}$ is the average energy liberated by CR ionization.
%This should \textbf{TODO:check} account for secondary ionizations as well.

% -----------------------------------
\subsubsection{Photoelectric heating}

\begin{figure}
    \includegraphics[width=\textwidth]{{fig/pe_heating}.pdf}
    \caption{
        PE heating expressions from BT94 and WD01
        Note, I rescaled WD01 expression
        from $10^{-26} \unit{erg\;s^{-1}}$
        to $10^{-24} \unit{erg\;s^{-1}}$
        in order to permit direct comparison of $\epsilon$ with BT94.
    }
    %\label{fig:peheat}
\end{figure}

Warning: notes on PE heating were for quick, one-off study, not meant for
long-term use.  Please clean up / check carefully before using.

$\Gamma_\mt{PE}$ is photoelectric heating; i.e., ``indirect'' radiation heating
from \texttt{PEFL\_VAR} liberating electrons on dust grains.  First, we need
the FUV flux:
\[
    G = \mathtt{PEFL\_VAR} + G_0 \, f_\mt{shield}\left(\rho,n,T\right) e^{-|z| / z_{0,FUV}} ,
\]
where $G$ and $G_0$ are normalized to the
Habing flux $1.6 \times 10^{-3} \unit{erg\;s^{-1}}$.
Josh adopts $G_0 = 1.69$ following Draine (1978).
Note: Josh's thesis defines $G_0$ as the dimension-ful Habing flux.
Here, $G_0 = \mathtt{he\_Gzero} = 1.69$ means a heating rate
$1.69 \times 1.6 \times 10^{-3} \unit{erg\;s^{-1}}$;
Bakes/Tielens (1994) use our convention, Weingartner/Draine use $G$ instead of
$G_0$.
%\textbf{TODO (20190619): I think Josh is trying to follow Weingartner/Draine convention.
%Change code variable name.}

The heating rate is from Weingartner and Draine
\href{http://adsabs.harvard.edu/abs/2001ApJS..134..263W}{2001ApJS..134..263W},
equation (44) and Table 2 with $R_V = 3.1$ and $b_c = 6 \times 10^{-5}$.
\begin{align*}
    G_\mt{f} &= G \sqrt{T} / n_\mt{e} \\
    \epsilon &= \frac{ 7.64 + 4.52 \times T^{0.132} }{ 1 + 0.04371 {G_\mt{f}}^{0.452} \left(1+0.00557{G_\mt{f}}^{0.675}\right) } \\
    \Gamma_\mt{PE} &= 10^{-26} \unit{erg\;s^{-1}} \times \epsilon G n / \rho
\end{align*}
The factor $n/\rho = 1/\mu$ gives a specific heating rate.
Note that $\epsilon$ influences how the radiation energy from $G$ gets divided between
thermal energy of gas versus dust; the heating rate for dust is
$\Gamma_\mt{PE\to dust} = 10^{-26} \unit{erg\;s^{-1}} \times (1-\epsilon) G n / \rho$.
whenever $\epsilon < 1$.

% fixed in june 2019 - but maybe good to review this 1 - epsilon calculation.

%For reasonable temperatures and densities, say
%$T < 10^4 \unit{K}$ and $n > 10^{-4} \unit{cm^{-3}}$,
%the factor $\epsilon > 1$ and so dust doesn't get any heating, all the energy goes
%into the liberated photo-electron.
%\textbf{TODO 2019 June 14} I'm not sure how to understand the $1-\epsilon$ in the
%code, one expects energy from FUV flux to be partitioned into dust and free
%electron (gas) heat.  Any energy dumped in dust will radiate away, but that
%ought to be accounted for in the dust temperature calculation.  Problem is
%$\epsilon$ does not lie within $[0, 1]$.
%\textbf{TODO 2019 June 19} I think this is resolved by rescaling from
%$10^{-26}$ to $10^{-24}$ in WD01 pre-factor.

%The Weingartner/Draine expression is a bit confusing - the range of values is
%quite large at the densities we are considering.  What is the range of
%validity?
%Also, Josh is adding a local column density approx, which counter-balances this
%effect somewhat.
In Torch, $n_e = x n_\mt{H}$ is taken from the ionization solver, which does
not track O, C, etc chemistry.  So our PE heating rates should be considered
fairly approximate.

Joung, M-MML (2006), Hill+ (2012), and Hill+ (2018) use a simpler form of this
expression:
\[
    \Gamma_\mt{PE} = 10^{-24} \unit{erg\;s^{-1}} \times \epsilon G n / \rho
\]
with $G = G_0 e^{-|z| / z_{0}}$ (this FUV scale height $z_0$ may differ from
grav, SN driving scale heights).
Joung, M-MML (2006) adopt $G_0 = 1.7$ and $\epsilon = 0.05$, yielding
$\Gamma_\mt{PE}/n = 0.85 \times 10^{-25} \unit{erg\;s^{-1}\;(H\,atom)^{-1}}$.
Hill+ (2018) adopt a larger fiducial value
$\Gamma_\mt{PE}/n = 3 \times 10^{-25} \unit{erg\;s^{-1}\;(H\,atom)^{-1}}$.

Kim and Ostriker (2017) do FUV from star particles.
They make two adjustments: (1) smooth linear turnoff towards high temperature,
using their mean molecular mass prescription (a proxy for ionization), and (2)
use a floor to represent ``metagalactic FUV'' that is present even in the
absence of star particles.  For cold gas with $\mu(T) = \mu_\mt{atomic}$, they
use $\Gamma = 0.0024 \times 2 \times 10^{-26} \unit{erg\;s^{-1}}$.

Josh also implemented (but now disabled) Wolfire 1995 and 2003 prescription.
Let's review these.

Wolfire+ 1995 equations (1--2) are:
\begin{align*}
    \Gamma &= 10^{-24} \unit{erg\;s^{-1}} \epsilon G_0 \\
    \epsilon
    &= \frac{0.049}{1 + \left( G_\mt{f} / 1925 \right)^{0.73}}
    + \frac{0.037 \left(T/10^4\unit{K}\right)^{0.7}}{1 + G_\mt{f} / 5000}
\end{align*}
where the efficiency $\epsilon$ is from Bakes \& Tielens 1994 and spans 0.001
to 0.1 for a range of $G_\mt{f}$ values.

Wolfire+ 2003 equations (19--20) are nearly the same, but adjusted to higher
PAH abundances:
\begin{align*}
    \Gamma &= 1.3 \times 10^{-24} \unit{erg\;s^{-1}} \epsilon G_0 \\
    \epsilon
    &= \frac{0.049}{1 + \left(G_\mt{f}/1925/\phi_\mt{PAH}\right)^{0.73}}
    + \frac{0.037 \left(T/10^4\unit{K}\right)^{0.7}}{1 + G_\mt{f}/5000/\phi_mt{PAH}}
\end{align*}
Wolfire+ 2003 adopt $\phi_\mt{PAH} = 0.5$ for their ``standard model''.

OTHER references: Wolfire 1995, Avillez and Breitschwerdt 2004

% -----
% a bit messy - discussion of this (1-epsilon) factor needs to be cleaned up
%
%% Question: how do we budget PEH heating to dust vs to gas?  I.e., how to explain
%% Josh's $(1-f)$ factor.
%%
%% ISRF: Weingartner/Draine quoted $5.33 \times 10^{-14} \unit{erg\;cm^{-3}}$ for
%% Habing (1968) and $8.93 \times 10^{-14} \unit{erg\;cm^{-3}}$ for Draine (1978).
%% Is it safe to assume equilibrium reaction?
%%
%% Goldsmith \href{https://ui.adsabs.harvard.edu/abs/2001ApJ...557..736G/abstract}{2001ApJ...557..736G}
%% and
%% Banerjee, Pudritz, Anderson
%% \href{http://adsabs.harvard.edu/abs/2006MNRAS.373.1091B}{2006MNRAS.373.1091B}
%% do not include PEH/external heating term, because FUV flux is shielded in
%% cores.
%%
%% Klessen/Glover 2014 review states dust heating in terms of $\Gamma_\mt{ext}$,
%% not separating contributions from CR and photo-heating.
%% The specific heating rate is of order
%% $\Gamma_\mt{ext} \sim 3 \times 10^{-24} \unit{erg\;s^{-1}} \, n/\rho$
%% which is like $10\times$ larger than the PE heating rate, so it's probably OK
%% to neglect the dust photo-heating that goes into gas rather than dust.
%%
%% If $\Gamma_\mt{ext}$ is dominated by photo-heating and not CRs, then  on the
%% basis of the cited heating rate, I think the use of $(1-\epsilon)$ in the code
%% can be adjusted by rescaling the Weingartner/Draine expression to
%% $10^{-24} \unit{erg\;s^{-1}}$, which lowers $\epsilon$ by $10$--$100$.
% -----

Krumholz's star formation textbook quotes a lower number:
\[
    \Gamma = 4 \times 10^{-26} G_0 \unit{erg\;s^{-1}}
\]
where I assume Habing is correct for ``local interstellar''.  For efficient
$\epsilon=0.05$, this implies total heating of gas and dust to be $8 \times
10^{-25} \unit{erg\;s^{-1}}$.
This is quite close to the value used by Joung and M-MML 2006, but somewhat
lower than the values considered by Hill+ (2018): $10^{-25}$ to $10^{-24}$
(pure PE gas heating, not the total FUV gas+dust heating).

% ----------------------------------
\subsubsection{Atomic cooling curve}

\begin{figure}
    \center
    \includegraphics[width=0.8\textwidth]{{fig/gas_atomic_cool}.png}
    \caption{
        Atomic cooling curve compared to Dalgarno/McCray 1972.
    }
    %\label{fig:peheat}
\end{figure}

$\Lambda_\mt{atom}$ is the cooling curve for atomic gas, originally used
in Joung and Mac Low 2006.  The references are:
Dalgarno and McCray \href{http://adsabs.harvard.edu/abs/1972ARA\%26A..10..375D}{1972ARA\&A..10..375D},
Sutherland and Dopita \href{https://ui.adsabs.harvard.edu/abs/1993ApJS...88..253S}{1993ApJS...88..253S}

Dalgarno/McCray include (1) free electron collisions that excite lines in both
ions and neutrals, (2) neutral hydrogen collisions in low-ionization regime to
excite fine structure lines of C${}^+$ and O.  They also state:
``The conversion of H to H${}_2$ substantially \emph{reduces} the cooling rate
of gas at low temperatures.  The spin change mechanism that [controls]
fine-structure excitation by neutral atomic hydrogen impact does not operate
with molecular hydrogen.'' They don't give rate estimates for H2 cooling, and
suggest just dividing atomic hydrogen rates by 4.

The validity of the atomic cooling curve at lower temperature is uncertain.

Comparison: Wolfire+ 1995 and 2003 explicitly write formulae for the
fine-structure line cooling from C II, O I, et cetera.

Dopita \& Sutherland is standard: recombination cooling and bremsstrahlung,
fewer physical processes need be considered in general.

% -------------------------------------
\subsubsection{Molecular cooling curve}

$\Lambda_\mt{mol}$ uses molecular cooling tables (\texttt{cool.dat}) adopted from
Neufeld+ 1993, Neufeld+ 1995; Banerjee papers 2004, 2006a, 2006b, 2015.

\textbf{Note:} H2 dissolution and formation timescales are comparable to MC
timescales, but we don't account for this chemistry!

Glover \& Clark
\href{https://ui.adsabs.harvard.edu/abs/2012MNRAS.421..116G/abstract}{2012MNRAS.421..116G}
have simulations with $n \sim 1$ to $100 \unit{cm^{-3}}$
with CO forming over $\abt 1 \unit{Myr}$.
Bigger simulation with colliding flows indicates $1$--$10 \unit{Myr}$,
Clark, Glover, Klessen, Bonnell
\href{https://ui.adsabs.harvard.edu/abs/2012MNRAS.424.2599C}{2012MNRAS.424.2599C}

Liszt
\href{https://ui.adsabs.harvard.edu/abs/2007A\%26A...461..205L/abstract}{2007A\&A...461..205L}
cites $\abt 10 \unit{Myr}$ in a neutral H I cloud of column density
$N(H) \sim 4 \times 10^{20} \unit{cm^{-2}}$ (which is low-ish compared to big
MW clouds, I think).

Useful overview of molecule cooling: Goldsmith and Langer
\href{https://ui.adsabs.harvard.edu/abs/1978ApJ...222..881G/abstract}{1978ApJ...222..881G}


% -------------------------------------
\subsubsection{Collisional gas-to-dust cooling}

$\Lambda_\mt{g\to d}$ cooling from gas colliding with dust, then dust
rapidly radiating away (dust acts as a heat sink).
\[
    \Gamma_\mt{g\leftarrow d}
    = \left(3.8 \times 10^{-33} \unit{erg\;cm^3\;s^{-1}} \right)
      \left(1 - 0.8 e^{-75/T_\mt{g}}\right)
      {T_\mt{g}}^{1/2} \left( T_\mt{g} - T_\mt{d} \right) {n_\mt{H}}^2 \rho^{-1}
    % this is volumetric rate erg /s /cm^3
    % convert to specific rate erg/s/g
\]
This is erg/s/g instead of erg/s/cm${}^3$ expressions quoted in dust heat/cool
calculation; they differ by a factor $\rho^{-1}$.


\begin{figure}
    \center
    \includegraphics[width=0.5\textwidth]{{fig/gas_cool}.pdf}
    \caption{
        Atomic cooling curve with interpolation at low temperature,
        and also dust cooling.
    }
    %\label{fig:peheat}
\end{figure}

% --------------------
% Dust Heating/Cooling
% --------------------
\subsection{Dust heat/cool}
obeys the equation
\begin{equation}
    n \frac{\dtl \epsilon}{\dtl t} =
    \Gamma_\mt{PE} + \Gamma_\mt{CR} + \Gamma_\mt{d\leftarrow g}
    - \Lambda_\mt{dust}
\end{equation}
where we have abruptly switched to volumetric heating rates, units:
$\mathrm{erg\;s^{-1}\;cm^{-3}}$.

\textbf{Method:} to solve for dust temperature, let $f\left(T_\mt{d}\right) = n \; \dtl\epsilon/\dtl t$.
Use Newton's method.  Iteratively compute:
\[
    T_\mt{d,new} = T_\mt{d} - \frac{ f\left(T_\mt{d}\right) }{ \frac{\ptl f}{\ptl T_\mt{d}} }
\]
until $T_\mt{d,new} = T_\mt{d}$ is converged.  This solution method, plus
commented out code relating to dust opacity, is described in the appendix of
Banerjee, Pudritz, Anderson \href{http://adsabs.harvard.edu/abs/2006MNRAS.373.1091B}{2006MNRAS.373.1091B}.
\textbf{TODO 2019 May 15: can we just use Brent's method instead?}

$\Gamma_\mt{PE}$ and $\Gamma_\mt{CR}$ are radiative and CR heating
respectively.  This replaces the commonly assumed ``external'' heating
$\Gamma_\mt{ext}$ from an assumed interstellar radiation field (ISRF).
\textbf{TODO 2019 May 14:} check that we reproduce correct ``external'' heating
rate with Habing FUV background?

$\Gamma_\mt{d\leftarrow g}$ is collisional heating of dust by gas from
Glover/Clark 2012, taken from Hollenbach/McKee.  Here $T_\mt{g}$ and $T_\mt{d}$
are gas and dust temperatures respectively.
\[
    \Gamma_\mt{d\leftarrow g}
    = 3.8 \times 10^{-33} \left(1 - 0.8 e^{-75/T_\mt{g}}\right)
    {T_\mt{g}}^{1/2} \left( T_\mt{g} - T_\mt{d} \right) {n_\mt{H}}^2
\]


$\Lambda_\mt{dust}$ is radiative dust cooling from Goldsmith 2001:
\[
    \Lambda_\mt{dust}
    =
    6.8 \times 10^{-33} {T_\mt{d}}^6 \; n\left(\mt{H}_2\right)
\]
in units $\mathrm{erg \; s^{-1} \; cm^{-3}}$.

This apparently disagrees with the result stated in Glover and Clark 2012
(reprinted in Klessen and Glover 2014 review).
I probably didn't account for or understand a necessary pre-factor.
Anyways, Glover and Clark (2012) state in equations (A7-8):
\[
    \Lambda_\mt{dust}
    =
    4\pi \rho_d \int_0^\infty B_\nu\left(T_\mt{d}\right) \; \kappa_\nu \dtl\nu
    \approx
    4.68 \times 10^{-31} {T_\mt{d}}^6 \; n
\]
in units $\mathrm{erg \; s^{-1} \; cm^{-3}}$.
Here $n$ is hydrogen number density, dust-to-gas mass ratio is probably $0.01$
(they claim it's a ``standard value'' but don't state the number),
and opacity $\kappa(\nu)$ is taken from Ossenkopf and Henning
\href{http://adsabs.harvard.edu/abs/1994A\%26A...291..943O}{1994A\&26A...291..943O}
for wavelengths $> 1 \unit{\mu m}$, and Mathis et al.~(1983) for shorter
wavelengths.  The Ossenkopf/Henning opacities are for non-coagulated grains
with thick ice mantles, intended for protostellar core modeling.

The use of molecular versus atomic number density can only explain a factor of
$2$ difference at most.  Why are these values so different?
Answer (2019 May 17): opacities, cannot be anywhere else.
MC conditions are quite different from protostellar core conditions, so
probably more physical to use Goldsmith, which indeed we use in Torch.

\begin{figure}
    \center
    \includegraphics[width=0.5\textwidth]{{fig/dust_opacity}.pdf}
    \caption{
        Compare dust opacities from Goldsmith versus Ossenkopf/Henning
    }
    %\label{fig:peheat}
\end{figure}

\textbf{Note (2019 May 15):} dust cooling equation is only good for
$T < 100 \unit{K}$ but we are using it all the way up to $T \sim 10^4
\unit{K}$.  In Draine book and other places, the dust cooling expression
plateaus.

Three papers cited for rates:
Goldsmith \href{https://ui.adsabs.harvard.edu/abs/2001ApJ...557..736G/abstract}{2001ApJ...557..736G},
Glover and Clark \href{https://ui.adsabs.harvard.edu/abs/2012MNRAS.421..116G/abstract}{2012MNRAS.421..116G},
Hollenbach and McKee \href{https://ui.adsabs.harvard.edu/abs/1989ApJ...342..306H/abstract}{1989ApJ...342..306H}.

\begin{figure}
    \center
    \includegraphics[width=0.85\textwidth]{{fig/dust_heat_cool_Tgas10}.pdf}
    \includegraphics[width=0.85\textwidth]{{fig/dust_heat_cool_Tgas10000}.pdf}
    \caption{
        Goldsmith versus Glover/Clark give slightly different dust
        temperatures.  Dust attains equilibrium temperature in a fraction of a
        second.
        Left column: net dust heating/cooling rate for given $T_\mt{dust}$ and
        $n_\mt{H}$.
        Right column: timescale for dust heating/cooling.
        Top four panels, $T_\mt{gas} = 10$; bottom, $T_\mt{gas} = 10^4$.
    }
    %\label{fig:peheat}
\end{figure}

% below units are only in flash4.2.2-rad/ , not part of main torch repo
% just various old versions of heating code
%
%\subsection{Banerjee/Seifried mol+dust module}
%
%Location: \\
%\texttt{
%HeatMain/phenHeat/RadHeat/molecules+dust
%}
%
%Code appears to harness FLASH's built-in Cooling unit?
%
%\subsection{JCIM?: phenHeat/IGWCurve}
%
%Location: \\
%\texttt{
%flash4.2.2-rad/source/physics/sourceTerms/Heat/phenHeat
%}
%
%I dunno why this was taken out of \texttt{HeatMain}.
%
%\subsection{Joung/Baczynski?: phenHeat with rad xfer}
%
%Location: \\
%\texttt{
%flash4.2.2-rad/source/physics/sourceTerms/Heat/HeatMain/phenHeat/RadHeat
%}
%
%What is \textbf{different}?
%
%\subsection{Joung/Baczynski?: phenHeat}
%
%Location: \\
%\texttt{
%flash4.2.2-rad/source/physics/sourceTerms/Heat/HeatMain/phenHeat
%}
%
%Heating: background FUV, user sets via \texttt{he\_peheat},
%and photo-electric heating via variable \texttt{PHHE\_VAR}.
%
%Cooling: Ryan Joung wrote a cooling function from some combination of Sarazin
%(1986), Raymond (1976), and Dalgarno \& McCray (1972).
%
%Sub-cycle as follows.  During $i$-th hydro step, compute the cooling time using
%an explicit method, where $n$ and $n+1$ index the sub-cycle steps within the
%$i$-th hydro step.
%\begin{align*}
%    \Delta t^n_\mathrm{sub} &= \min\left(
%      f_\mathrm{sub} t^{n}_\mathrm{cool},\,
%      t_{\mathrm{hy},i+1} - t^n_\mathrm{sub}
%    \right) \\
%    t_\mathrm{cool} &= \frac{\epsilon^{n}}{\left| \mathrm{d}\epsilon / \mathrm{d}t \right|}
%    = \frac{\epsilon^{n}}{\left| n^2 \mathcal{L}/\rho \right|}
%    = \frac{\epsilon^{n}}{ \left| n \Gamma(\vec{r}) - n^2 \Lambda(T^n) \right| /\rho }
%\end{align*}
%
%With subcycle timestep in hand, advance subcycle energy to $\epsilon^{n+1}$
%with an explicit method:
%\[
%    \frac{\epsilon^{n+1} - \epsilon^{n}}{\Delta t^n_\mathrm{sub}}
%    = n^2\mathcal{L}/\rho = \left[ n\Gamma\left(\vec{r}\right) - n^2 \Lambda(T^n)\right]/\rho
%\]

% =============
% Stellar winds
% =============
\section{Stellar winds}

Implemented in: \\
\texttt{
source/Particles/ParticlesMain/active/Sink/Couple\_AMUSE/wind
}

Star particles have the following important properties:
\begin{verbatim}
    BGDY_PART_PROP              background gas density at time of wind start
    CREATION_TIME_PART_PROP     creation time
    DMDT_PART_PROP              star's mass loss rate
    POSX_PART_PROP              x position
    POSY_PART_PROP              y position
    POSZ_PART_PROP              z position
    VELW_PART_PROP              star's mass loss rate
\end{verbatim}
which are used to inject winds.

Note: \texttt{ref\_radius} was formerly used in particle mark refine/derefine.
But particle mark refine/derefine criterion has been moved to junk, so not used?.
%\textbf{TODO (2019 Jul 03) Why was the prtl refinement removed? Or was it
%refinement implemented elsewhere?}

There is an adjustment to the injection scheme for winds without mass-loading.
Each cell is weighted by the solid angle $\Omega$ computed from
\[
    \cos^2 \left(\Omega/4\right)
    = \frac {1 + \frac{1}{2} \frac{\Delta^2}{r^2}}
            {1 + \frac{1}{2} \frac{\Delta^2}{r^2} + \frac{\Delta^4}{r^4}}
\]
which should make the mass density decrease as $r^{-2}$ (i.e., constant mass in
each shell, as expected for free expansion), according to commit
flash4.2.2-rad \texttt{88fc47a}.

Note, the feature for adding thermal energy to momentum-conserving wind is
relatively new.  Josh has some comments to this effect,
correct amount of thermal energy to add is not 100\% clear.
Commit is FLASH4.5 \texttt{00272ec}.


% ========================
% Random Supernova Driving
% ========================
\section{Supernova Driving}

% ---------------------------------------
% Josh's implementation of Simpson method
% ---------------------------------------
\subsection{Wall: momentum+thermal SN driving}

\texttt{
source/Particles/ParticlesMain/energyInjection
}

% ----------------------------------------
% Joung/M-MML/Baczynski thermal SN driving
% ----------------------------------------
\subsection{Joung: pure thermal SN driving}

Principle: deposit 100\% thermal energy in a sphere.
Mass of 60 Msun is chosen to prevent unphysical cooling (either at too hot or
too cold temperatures).

%\textbf{TODO 2019 May 16 - understand the choice of mass shells, what's going
%on}

References: Joung \& Mac Low (2006); Joung, Mac Low, \& Bryan (2009);
Gatto et al.~(?) (2015); Ib\'{a}\~{n}ez-Mej\'{i}a et al.~(2016,2017).

Key units are: \\
\texttt{
source/physics/sourceTerms/SN \\
source/physics/sourceTerms/SN/SNSBrandom \\
source/physics/sourceTerms/SN/SNSBrandom/Clustersink \\
source/physics/sourceTerms/SN/SNSBrandom/SNsink \\
source/physics/sourceTerms/SN/SNSBrandom/SBpeak \\
source/physics/sourceTerms/SN/SNSBrandom/SNejecta \\
source/physics/sourceTerms/SN/SNSBrandom/SNkinetic \\
source/Particles/ParticlesMain/sb \\
source/Particles/ParticlesMain/active/Sink/Clustersink \\
source/Particles/ParticlesMain/active/Sink/SNsink \\
}

Further units, used with SNSB unit and Juan's stratboxes: \\
\texttt{
source/Particles/ParticlesMain/passive/RungeKuttaThermo \\
source/Particles/ParticlesMain/passive/MonteCarlo \\
source/Particles/ParticlesMapping/NoMap \\
source/Particles/ParticlesMapping/InZone \\
source/Particles/ParticlesInitialization/NParPerZone \\
source/Particles/ParticlesInitialization/NParPerZoneMC
}

To use, request \texttt{sourceTerms/SN/SNSBrandom} and/or one of its sub-units,
and the necessary particles unit will be loaded.

A useful problem to test this unit on is
\texttt{
source/Simulation/SimulationMain/SNSBbox \\
source/Simulation/SimulationMain/SNSBbox\_AT
}

Aaron ported from FLASH4.2.2 to FLASH4.5 on 2019 May 01.  The code is deeply
interwoven into vanilla FLASH source, but this is probably the cleanest way to
do it (avoids copy-paste, and allows the various init/mapping units to be used
in the most modular way).

Easiest way to review changes is by diffing prtl directory directly (vanilla
FLASH vs Torch), but here is a guide to help disentangle the distinct changes
since there's a lot going on.

\textbf{TODO May 2019}: only some not all prtl functionality tested.

\begin{itemize}

    \item SN unit is hardcoded into Driver sourceTerms, initSourceTerms, computeDt,
          finalizeSourceTerms.

      \item \texttt{Particles/Particles.h} edited to register new particle
          types, mapping, and initialization units. Note that
          \texttt{RungeKuttaStatic} is hidden under
          \texttt{ParticlesMain/sb/}.

    \item \texttt{ParticlesMain/passive/Config} edited to register new passive
        integrators RungeKuttaThermo, MonteCarlo.

    \item New \textbf{ParticlesMain/passive} particle types \texttt{RungeKuttaThermo},
        \texttt{MonteCarlo}, and \texttt{sb/RungeKuttaStatic} require
        registration.

        Define new \texttt{Particles/Particles.h} constants
        \texttt{RUNGEKUTTASTATIC}, \texttt{RUNGEKUTTATHERMO}, and
        \texttt{MONTECARLO}.

        Add stubs: \\
        \texttt{Particles/localAPI/pt\_advanceRKstatic.F90 \\
                Particles/localAPI/pt\_advanceRKthermo.F90 \\
                Particles/localAPI/pt\_advanceMonteCarlo.F90
        } \\
        and register them in \texttt{Particles/localAPI/Makefile} so that code
        knows how to advance prtl from different sub-units.

        Edit: \\
        \texttt{
            Particles/ParticlesMain/Particles\_advance.F90
        }
        to add new case logic for prtl advance.

        Finally, for \texttt{RUNGEKUTTATHERMO} alone, edit: \\
        \texttt{
            Particles/ParticlesMain/Particles\_data.F90 \\
            Particles/ParticlesInitialization/pt\_setDataStructures.F90
        } \\
        to get

        \textbf{TODO/BEWARE!} \texttt{Particles\_advance} is overwritten by
        AMUSE coupling code... this will probably break passive tracers.


    \item \textbf{ParticlesInitialization} sub-units \texttt{NParPerZone} and
        \texttt{NParPerZoneMC} require registration.

        Define new \texttt{Particles/Particles.h} constants
        \texttt{NPPARPERZONE} and \texttt{NPPARPERZONEMC}

        Add stubs: \\
        \texttt{Particles/localAPI/pt\_initPositionsNParPerZone.F90 \\
                Particles/localAPI/pt\_initPositionsNParPerZoneMC.F90
        } \\
        and register them in \texttt{Particles/localAPI/Makefile} so that code
        knows how to initialize prtl from different sub-units.

        Edit: \\
        \texttt{
            Particles/ParticlesInitialization/Particles\_init.F90 \\
            Particles/ParticlesInitialization/Particles\_initPositions.F90 \\
            Particles/ParticlesMain/Particles\_data.F90 \\
        }
        to add new case logic for prtl init, and to add some config options.

        QUESTION: can we implement as "CUSTOM" instead of named sub-units?
        ANSWER: maybe.  But, this looks possibly useful for things other than
        SN unit.

    \item \textbf{ParticlesMapping} sub-units \texttt{InZone} and
        \texttt{NoMap} require registration.

        Define new \texttt{Particles/Particles.h} constants
        \texttt{INZONE} and \texttt{NOMAP}

        Add stubs: \\
        \texttt{Particles/localAPI/pt\_mapFromMeshInZone.F90 \\
                Particles/localAPI/pt\_mapFromMeshNone.F90
        } \\
        and register them in \texttt{Particles/localAPI/Makefile} so that code
        knows how to map prtl from different sub-units.
        \textbf{WARNING} for ParticlesMapping ALSO register in
        \texttt{pt\_interface.F90}, I don't understand why the init/integrate
        methods don't need to be registered here for subunits to find........

        Edit: \\
        \texttt{
            Particles/ParticlesMapping/Particles\_mapFromMesh.F90
        } \\
        to add new case logic for prtl mapping.

    \item Why do we have to register all this stuff in localAPI?
        Answer: PRTL code has a bunch of sub-units.
        Even if our simulation doesn't have all sub-units, it must know the
        methods from all sub-units in order to compile.

\end{itemize}
%The hardcoding into FLASH source makes future Torch upgrades more difficult.
%However, the edits are small and alternatives had their own cons.

Joung/CB SN heating is applied as a heating rate via \texttt{PHHE\_VAR}, which
does not play nice with the rad xfer (fluxes get cleared in subcycles, but the
SN heating rate needs to be applied continuously).

Tentatively, I (AT) have reverted to older direct \texttt{EINT\_VAR} and
\texttt{ENER\_VAR} injection method, which redistributes the mass within the
injection sphere.  \textbf{TODO 20190614} this also breaks momentum
conservation, to be addressed by AT.

To use SB particle unit, currently you have to declare particle type
\texttt{sb} in your problem Config.  It appears to work fine with Josh's code.
Although, need to make sure that particles are requested/addressed by name
explicitly everywhere, and does not assume / require only two prtl types in
logic.

Not sure how to make Particle unit advance work nicely with Josh's setup.
In flash4.2.2-rad, CB manually registered  \texttt{Particles\_advance} method...
commit from like 2013...

% --------------------------------------------------------
% Joung/M-MML/Baczynski thermal SN driving, particle stuff
% --------------------------------------------------------
%\subsection{Particle stuff in Joung unit}



% --------------------------------
% Comparison of SN driving schemes
% --------------------------------
\subsection{Compare to other SN driving schemes}

Key point: if you don't resolve the radius where shell formation occurs with at
least 3--5 cells, you will not get energy/momentum injection correct.  All
authors broadly agree on this.

Remark (M-MML): evaporation of shell mass into hot phase might matter and is generally
neglected.

Martizzi+ \href{https://ui.adsabs.harvard.edu/abs/2016MNRAS.459.2311M/abstract}{2016MNRAS.459.2311M}
SN feedback in local stratified ISM for MW- and ULIRG-like disks
using the subgrid model from Martizzi+
\href{https://ui.adsabs.harvard.edu/abs/2015MNRAS.450..504M/abstract}{2015MNRAS.450..504M}
\begin{itemize}
    \item Procedure: inject uniform thermal energy and radial momentum into a
        sphere of radius $2 \Delta x$.  Use Martizzi+ 2015 equations (9-10) to
        figure out how much to inject, for forward shock $R = R_\mt{inj}$:
\[
    E_\mt{th}\left(R, n_\mt{H}, Z\right) = E_\mt{th,0}
    \begin{cases}
        1                                       &, R < R_\mt{c} \\
        \left(R / R_\mt{c}\right)^\alpha        &, R_\mt{c} < R < R_\mt{r} \\
        \left(R_\mt{r} / R_\mt{c}\right)^\alpha &, R_\mt{r} < R
    \end{cases}
\]
        where $R_\mt{c}$ is cooling radius where snow-plough shell forms, and
        $R_\mt{r}$ is late-time radius where thermal energy becomes roughly
        constant.  For momentum (changing Martizzi's $P_r$ to $p_r$) use:
\[
    p_\mt{r}\left(R, n_\mt{H}, Z\right) = P_\mt{0}
    \begin{cases}
        \left(R / R_0\right)^{1.5}      &, R < R_\mt{b} \\
        \left(R_\mt{b}/R_0\right)^{1.5} &, R > R_\mt{b}
    \end{cases}
\]
        where $R_0$ is scale radius used to extrapolate S-T power law to small
        radii, and $R_\mt{b}$ is where SNR momentum becomes roughly constant.
        The hierarchy is $R_0 < R_\mt{c} < R_\mt{b} < R_\mt{r}$.
        S-T scale $<$  cooling radius $<$ radius where momentum stabilizes $<$
        radius where thermal energy stabilizes.  See Figure 4 of Martizzi.

        Compute mean density, metallicity within injection sphere.
        Use Martizzi+ 2015 paper fitting functions, equation (11), to determine
        $\alpha$, $R_\mt{c}$, $R_\mt{r}$, $R_0$, $R_\mt{b}$.
        Then, use equations (9-10) reproduced above to inject appropriate
        energy and momentum.

    \item They sampled densities $0.01$, $1$, $100$, and $10^4 \unit{cm^{-3}}$.
        with solar metallicity $Z_\sun$.

    \item When SNR cooling radius is well resolved ($R_\mt{inj} < R_c$)
        the model ``reduces to injecting $10^{51} \unit{erg}$ of thermal energy
        per SN and letting the simulation code explicitly capture the
        conversion into radial momentum''.

    \item Note that cooling radius $R_\mt{c}$ is...
\[
    R_\mt{c} \approx 14 \unit{pc}
        \; \left(n_\mt{H}/1\unit{cm^{-3}}\right)^{-3/7}
        \; \left(E_\mt{tot}/10^{51}\unit{erg}\right)^{2/7}
        \; \left(Z/Z_\sun\right)^{-1/7}
\]
    based on Cioffi+ (1988) and Thornton+ (1998).  In the high-resolution sims,
    Martizzi resolve $R_\mt{c}$ with at least 50 cells.
\end{itemize}


Miao Li+ (incl Jerry, Greg)
\href{https://ui.adsabs.harvard.edu/abs/2015ApJ...814....4L/abstract}{2015ApJ...814....4L}
also simulated SNe in uniform and inhomogeneous environments.
See Figure 2 for energy, momentum evolution with radius.
However, they don't provide any sub-grid model fitting functions.

Gatto+ \href{https://ui.adsabs.harvard.edu/abs/2015MNRAS.449.1057G/abstract}{2015MNRAS.449.1057G}
choose $R_\mt{inj}$ to contain $10^3 M_\sun$, and require
$R_\mt{inj} \geq 4 \Delta x$ (so, in dense environment.  Procedure is as
follows...
\begin{itemize}
    \item If SN is resolved, dump $10^{51} \unit{erg}$ pure thermal energy.
        Gas is heated to $8 \times 10^6 \unit{K}$ ($\mu = 1.3 m_\mt{P}$, H and
        He, ignore lecs) or $3 \times 10^6 \unit{K}$ ($\mu =  0.5 m_\mt{P}$, H
        and lec plasma, ignore He).
    \item If SN is unresolved, $R_\mt{inj}$ contains too much mass and would
        over-cool.  Then, inject momentum instead:
\[
    p_\mt{ST} = 2.6\times10^5 \unit{M_\sun\;km\;s^{-1}}
        \left(E/10^{51}\unit{erg}\right)^{16/17}
        \left(\bar{n}/1\unit{cm^{-3}}\right)^{-2/17}
\]
    Flatten density within injection region, add velocity, and heat the
    injection region to $10^4\unit{K}$.
\item They note that result should under-estimate momentum input;
    during snowplough phase, SNR may continue to gain momentum.
    Also, lacks strong shocks and hot gas.
\end{itemize}

Simpson, Bryan, Hummels, JPO
\href{https://ui.adsabs.harvard.edu/abs/2015ApJ...809...69S}{2015ApJ...809...69S}
\begin{itemize}
    \item Injection stencil is 3x3x3, so $r_\mt{inj} = 1.5 \Delta x$ and will
        put energy into 4x4x4 cells.
    \item If $4.5 \Delta x < R_\mt{PDS}$ (same as Martizzi $R_\mt{C}$),
        then just dump thermal energy.  Otherwise, do:
\[
    f_\mt{kin} = 3.97 \times 10^{-6} \mu n_0 {R_\mt{PDS}}^7 {t_\mt{PDS}}^{-2}
        \Delta x^{-2} {E_\mt{51}}^{-1}
\]
    \item To evaluate $t_\mt{PDS}$ and $r_\mt{PDS}$ use equations (17-18),
        which also encode density and metallicity dependence.
        For solar metallicity,
        $R_\mt{PDS} = 18.5 \unit{pc} \; {E_{51}}^{2/7} {n_0}^{-3/7}$
        is quite similar to expression of Martizzi
    \item Density range: 0.5 to $500 \unit{cm^{-3}}$, they show test of their
        prescription for $n = 49 \unit{cm^{-3}}$.
\end{itemize}

Kim \& Ostriker
\href{https://ui.adsabs.harvard.edu/abs/2015ApJ...802...99K}{2015ApJ...802...99K}
\begin{itemize}
    \item Inject a small region, $r_\mt{init} = \Delta x$ and so touch 2x2x2 cells (8 zones)
    \item $r_\mt{sf} = 22.6 \unit{pc} \; {E_{51}}^{0.29} {n_0}^{-0.42}$,
        note that $2/7 = 0.286$ and $-3/7 = -0.428$ so basically the same as
        all the other authors
    \item \textbf{Particularly useful for discussing $r_\mt{init}$ convergence,
        which I haven't seen other papers explore}
    \item Density range: 0.1 to $100 \unit{cm^{-3}}$
\end{itemize}

Kim \& Ostriker TIGRESS
\href{https://ui.adsabs.harvard.edu/abs/2017ApJ...846..133K/abstract}{2017ApJ...846..133K}
use a complex prescription with ST, EJ, MC regimes.
\begin{enumerate}
    \item Look at spheres from $R_\mt{min} = 3\Delta x$ to $R_\mt{max}$,
        varying $R$ by $\Delta x/2$.  For each sphere, compute $M/M_\mt{sf}$
        where $M_\mt{sf} = 1679 M_\sun \left(n_\mt{H}/1\unit{cm^{-3}}\right)^{-0.26}$
    \item IF $M > M_\mt{sf}$ at $R_\mt{min}$, then ST is unresolved.
        Inject momentum onto grid, using prescription from hi-res sims of SNRs
        in two-phase medium.  No thermal energy, assume all cooled and gone(?).
    \item IF $0.027 M_\mt{sf} < M < M_\mt{sf}$ at $R_\mt{min}$, then ST is
        resolved.  Inject 72\% thermal and 28\% kinetic energy onto grid.
    \item IF $M < 0.027 M_\mt{sf}$ at $R_\mt{min}$, increase $R$ until we reach
        $M \geq 0.027 M_\mt{sf}$.  Then do ST injection.
    \item IF $M < 0.027 M_\mt{sf}$ at $R_\mt{max}$ (can't find enough mass),
        then ejectra in free expansion (EJ regime).  Dump pure KE.
\end{enumerate}
They have some comments:
\begin{itemize}
    \item For ST regime, dumping thermal energy + redistribute mass is supposed
        to produce same results.  However, they claim you must re-distribute
        mass within region to avoid overcooling, for $M$ close to $M_\mt{sf}$.
    \item EJ, ST regimes: redistribute mass/velocity/energy.  For MC regime,
        do not reset gas properties, just add momentum to all zones.
    \item Argue Gatto+ 2015 feedback scheme is ``marginal'', and that they
        needed to redistribute mass in injection sphere, and so they got too
        high SFR.
    \item Figure 6: the prescription looks kind of bonkers at high density,
        artifact of fixed (low) spatial resolution.
        This is the MC regime; discard thermal energy, just keep momentum.
\end{itemize}

Hopkins FIRE
\href{https://ui.adsabs.harvard.edu/abs/2014MNRAS.445..581H/abstract}{2014MNRAS.445..581H}
\begin{itemize}
    \item Take
        $R_\mt{cool} \approx 28 \unit{pc} \; {E_{51}}^{0.29}
            \left(n / 1 \unit{cm^{-3}}\right)^{-0.43}
            \left(Z/Z_\sun + 0.01\right)^{-0.18}$.
    \item When S-T phase is resolved, deposit correct amounts of thermal
        and kinetic energy (momentum) directly.
    \item When not, deposit momentum boosted by $\abt {M_\mt{swept-up}}^{0.5}$,
        depending on how much mass got swept up by hot gas in the S-T stage.
\[
    \vec{p} = \vec{p}_\mt{ej} \sqrt{ \frac{M_\mt{ej} + M_\mt{enc}(<R_\mt{cool})}{M_\mt{ej}} }
\]
    and also decrease the thermal energy,
\[
    E_\mt{thermal,shocked} \propto \left(R/R_\mt{cool}\right)^{-6.5}
\]
    which reduces its ability to do any more work.
\end{itemize}

% -----------------------------------------
% Comparison of superbubble driving schemes
% -----------------------------------------
\subsection{Superbubble schemes and behavior}

Kim Ostriker Ralaenu

Useful, recent discussion of sub-grid SNe schemes and winds by Chia-Yu Hu, see
Section 4.3:
\href{https://ui.adsabs.harvard.edu/abs/2019MNRAS.483.3363H/abstract}{2019MNRAS.483.3363H}


% ================
% Code development
% ================
\section{Code development}

\subsection{Generalities}

Google's code style guides, more for philosophy insight rather than any
specific actionables:
\begin{itemize}
    \item Python style: \url{https://google.github.io/styleguide/pyguide.html}
    \item Documentation style: \url{https://developers.google.com/style/}
\end{itemize}
See in particular the ``Parting Words'' in the Python style guide.

``Maintenance cost is quadratic in number of features'':
\url{https://statmodeling.stat.columbia.edu/2019/05/10/maintenance-cost-is-quadratic-in-the-number-of-features/#comments}
To mitigate -- be clear about what use cases are supported or not -- reduce
maintenance cost to linear (though for a tower of features, that also means
code works only for one use case).

\subsection{FLASH add-ons/edits}

Try to follow the FLASH idiom.  Read (and maybe re-read) Chapter 4 of the FLASH
manual before making large code changes.  Tips on FLASH architecture: \\
\url{http://flash.uchicago.edu/site/flashcode/user_support/tipsDoc.py?submit=tips_arch.txt}

Useful FLASH tutorial presentations:
\begin{itemize}
    \item FLASH center 2012: \url{http://flash.uchicago.edu/site/flashcode/user_support/tutorial_talks/home.py?submit=RAL_May2012.txt}
    \item collective wisdom -- sfocu may be worth investigating for automated testing
        \url{http://flash.uchicago.edu/site/flashcode/user_support/tutorial_talks/June2009/CollectiveWisdom.pdf}
\end{itemize}

To overwrite a FLASH-default file, you can put stuff in
\texttt{Simulation/SimulationMain/YourProblem}.
%Try to avoid putting in FLASH source if you can, it makes things messy.

Idea: when you add a new FLASH sub-unit, it's easy to copy-paste a file and
just add stuff to it.  But, it can be hard to maintain because you're
duplicating code:
\begin{verbatim}
    Foo/
        Foo.F90
        Foo_bar.F90
        Foo_interface.F90
        FooMain/
            Foo.F90
            Foo_bar.F90
            FooLoremIpsum/
                Foo.F90      # contains lots of duplicate code
                Foo_bar.F90
\end{verbatim}

Consider trying:
\begin{verbatim}
    Foo/
        Foo.F90
        Foo_bar.F90
        Foo_lorem.F90
        Foo_ipsum.F90
        Foo_interface.F90
        FooMain/
            Foo.F90          # modify to allow for LoremIpsum features
            Foo_bar.F90
            FooLoremIpsum/
                Foo_lorem.F90  # no duplication!
                Foo_ipsum.F90
\end{verbatim}


Config file directives:
\url{http://flash.uchicago.edu/site/flashcode/user_support/flash4_ug_4p5/node35.html#SECTION03252000000000000000}

PARAMESH manual is no longer posted online, but available on the Wayback
machine.  Last snapshot was 2016 April 21.
Visit
\url{https://web.archive.org/web/*/http://www.physics.drexel.edu/~olson/paramesh-doc/Users_manual/amr.html}

Some thoughts on community code development:
\url{https://extremecomputingtraining.anl.gov//files/2015/03/commcodesAnshu-aug8-830.pdf}
\url{http://flash.uchicago.edu/cc2012/}

The FLASH center uses 3 space indents.  No tabs, 2 space continuations(?).
JW/JCIM/CB/MKRJ seem to use some mix of 2 space, 4 space, and tab.

Git forking workflow: in your forked repo, do NOT press the SYNC button on
bitbucket to grab upstream commits.
It generates an extra merge commit, cluttering history.
Maybe there's a workaround and I haven't discovered it.
It's easier to do \texttt{git merge upstream/master};
so long as you didn't touch master (don't touch master!!) it will do a
fast-forward merge and life will be good.


% -------------
% Pull Requests
% -------------
\subsection{Pull Requests}

To submit a pull request, you'll generally want to make a branch, do some work,
then push to your local repo.

We are following the Git forking workflow, see:
\url{https://www.atlassian.com/git/tutorials/comparing-workflows/forking-workflow}
and
\url{https://gist.github.com/Chaser324/ce0505fbed06b947d962}

If possible, try to rebase to most recent commit before you submit PR.

To clean now-removed remote branch pointers, do \texttt{git remote prune origin}.


% =====================================
% Lessons Learned, Cheatsheets, and FAQ
% =====================================
\section{Lessons Learned, Cheatsheets, and FAQ}

% ------------------------
% FLASH runtime parameters
% ------------------------
\subsection{FLASH Misc Gotchas}

Warning: sanitizing doesn't do anything apparently?
All the fixing up comes from floors (smlrho, smallp, etc).
Sources: Josh said so, and also
\url{http://flash.uchicago.edu/pipermail/flash-users/2009-May/000467.html}

Some snags with USM and paramesh
\begin{itemize}
    \item checkerboard B field structure with USM:
        \url{http://flash.uchicago.edu/pipermail/flash-users/2017-October/002427.html}

    \item spurious waves at refinement boundaries:
        \url{http://flash.uchicago.edu/pipermail/flash-users/2016-May/001956.html}

    \item more discussion about refinement boundary noise with usm...
        \url{http://flash.uchicago.edu/pipermail/flash-users/2019-March/002823.html}
\end{itemize}

%Important change: in the file
%\texttt{Grid/GridMain/paramesh/interpolation/Paramesh4/prolong/Config},
%comment out
%\texttt{PPDEFINE PM\_OPTIMIZE\_MORTONBND\_FETCHLIST}
%to prevent weird segfaults.  See the FLASH-USERS thread:
%\url{http://flash.uchicago.edu/pipermail/flash-users/2019-May/002901.html}


% ------------------------
% FLASH runtime parameters
% ------------------------
%\subsection{FLASH runtime parameters}


% ---------------------
% FLASH ifdef variables
% ---------------------
\subsection{FLASH ifdef variables}

\texttt{PE\_HEAT} indicates whether code tracks FUV photons or not; A.~Tran
changed behavior from toggling PE heating in general to only toggling FUV
photon heating from stars, because we need to be able to put in background PE
heating separate from stars as well.


% ---------------
% FLASH variables
% ---------------
\subsection{FLASH variable cheat sheet}

\subsubsection{FLASH field variables}

First grouped by unit, then ordered alphabetically?

\begin{verbatim}
PHIO_VAR    photoionization rate [1/s],                      ParticlesMain/ray_pe/MultiSourceSimple
PHHE_VAR    photoionization heating rate [erg/s],            ParticlesMain/ray_pe/MultiSourceSimple
PEFL_VAR    photoelectric flux [Habing units, scaled erg/s], ParticlesMain/ray_pe/MultiSourceSimple
UVFL_VAR    scratch var for photoionizing flux [erg/s],      ParticlesMain/ray_pe/MultiSourceSimple
FUFL_VAR    scratch var for photoelectric flux [erg/s],      ParticlesMain/ray_pe/MultiSourceSimple
\end{verbatim}

Gravity unit and Torch star-(sink,gas) bridge variables

\begin{verbatim}
BGPT_VAR    grav potential of star prtl, from Poisson eqn
BGAX_VAR    grav accel due to star prtl, x-component
BGAY_VAR    grav accel due to star prtl, y-component
BGAZ_VAR    grav accel due to star prtl, z-component
DENS_VAR    density of gas
GPOT_VAR    grav potential of gas (DENS_VAR)
PDEN_VAR    density of star prtl, mapped to mesh
SGAX_VAR    grav accel due to sinks, x-component
SGAY_VAR    grav accel due to sinks, y-component
SGAZ_VAR    grav accel due to sinks, z-component
\end{verbatim}

\subsubsection{FLASH particle variables}

Superbubble particles introduce several properties:

\begin{verbatim}
NSN_PART_PROP       number of SN in SB              ParticlesMain/sb
TCRT_PART_PROP      creation time of SB             ParticlesMain/sb
TSN_PART_PROP       time interval between SNe       ParticlesMain/sb
SBID_PART_PROP      SB unique identifier            ParticlesMain/sb
\end{verbatim}

In the supernova unit code, tracer particles overload some particle properties
in order to grab information about the surrounding gas.
The mesh-to-particle map happens in
\texttt{ParticlesMapping/Particles\_mapFromMesh.F90} ; the overloaded
properties are:

\begin{verbatim}
NSN_PART_PROP       maps to grid's TEMP_VAR        ParticlesInitialization/pt_setDataStructures
MASS_PART_PROP      maps to grid's DENS_VAR        ParticlesInitialization/pt_setDataStructures
\end{verbatim}

OK stuff for radiation and amuse coupling

\begin{verbatim}
EION_PART_PROP      energy of ionizing photons
NION_PART_PROP      number of ionizing photons
SIGH_PART_PROP      gas cross-section to ionizing radiation
EPEP_PART_PROP      Average energy of PE photon
NPEP_PART_PROP      number of PE photons
SIGD_PART_PROP      dust cross-section to PE photons
\end{verbatim}

Stuff for winds coupled to AMUSE

\begin{verbatim}
DMDT_PART_PROP      wind mass loss rate
VELW_PART_PROP      wind mass loss rate
\end{verbatim}

%\subsubsection{FLASH species variables}

% ---------------
% Frequently asked questions
% ---------------
\subsection{Frequently asked questions}

\begin{itemize}

    \item \textbf{Question:} How long does a run take?

    \textbf{Answer:} on Columbia University's cluster ``Habanero'' (Intel
    E5-2650v4, 2.2GHz; 5.3 GB/core; 24 cores/node; EDR Infiniband), with 6
    refinement levels, block size $16^3$, and 186 cores, evolving JUST GAS (no
    N-body, no feedback) to $3 \unit{Myr}$ took $10 \unit{days}$ wall clock
    time or $\abt 50,000$ CPU-hours.
    This was a very crude test run, not sure if I used bridge code or plain
    FLASH.  Not using parallel HDF5.
    A production run will take $100,000$s of CPU-hours.

    \item \textbf{Question:} How do I add a new runtime parameter?

    \textbf{Answer:}
        \begin{enumerate}
            \item define in \texttt{Simulation\_data.F90}
            \item define and document in \texttt{Config}
            \item in \texttt{Simulation\_init.F90}, call \texttt{RuntimeParameters\_get(...)}
            \item add to your \texttt{flash.par}
        \end{enumerate}

    \item \textbf{Question:} How do I track down a variable or error message?

    \textbf{Answer:} in your object directory, do something like
        \texttt{grep -rinT -I "INSERT ERROR MESSAGE" *}

    \item \textbf{Question:} How do I decipher FLASH units, runtime parameters,
        and setup calls?

    \textbf{Answer:} look at
        \texttt{bin/setup\_shortcuts.txt},
        \texttt{object/setup\_call},
        \texttt{object/setup\_params},
        \texttt{object/setup\_units}.

    \item \textbf{Question:} How much memory does FLASH use?

    \textbf{Answer:} estimate as:
    \begin{align*}
        8 \unit{bytes}
        &\times \left( \mathtt{NUNK\_VARS} + \mathtt{NFACEVAR} + \mathtt{NFLUXVAR} \right) \\
        &\times \left( \mathtt{NXB} + 2\times\mathtt{GUARDCELLS} \right)
        \left( \mathtt{NYB} + 2\times\mathtt{GUARDCELLS} \right)
        \left( \mathtt{NZB} + 2\times\mathtt{GUARDCELLS} \right) \\
        &\times \mathtt{MAXBLOCKS}
    \end{align*}
    These variables are defined by \texttt{Flash.h}, in the object directory.
    For the simulation \texttt{Cube} with unsplit MHD solver,
    $16^3$ blocks, and \texttt{MAXBLOCKS=300}, the memory usage is:
    \[
        8 \unit{bytes} \times (46 + 2 + 12) \times (16 + 12)^3 \times 300
        \approx 3.2 \unit{GB}
    \]

    \item \textbf{Question:} with FLASH4.2.2, how do I compile with large
        values of MAXBLOCKS?  Setup works, but make fails.

    \textbf{Answer:} in \texttt{Makefile.h}, request the compiler flag
        \texttt{-mcmodel=medium}.  Search flash-users mailing list for
        ``mcmodel'' to learn more.

    \item \textbf{Question:} ``needed full blk... but could not find it or only
        found part of it in the message buffer. Contact PARAMESH developers for
        help.''

    \textbf{Answer:}
    \url{http://flash.uchicago.edu/pipermail/flash-users/2015-February/001637.html}
    See also
    \url{http://flash.uchicago.edu/pipermail/flash-users/2017-April/002262.html}

    \item \textbf{Question:} when restarting, what numbers should I use for
        checkpoints and plotfiles?

    \textbf{Answer:} checkpoints, use the files you want to restart from.
        Plot files, use the next plot file numbers to create.

    \item \textbf{Question:} FLASH performance?

    \textbf{Answer:} Mordecai's proposal quotes 150e-6 core-seconds per zone
    update.  I (AT) measured very similar on Habanero cluster at Columbia,
    which is Intel Xeon E5-2650v4.

    Sean Couch reports on FLASH users list (November 2014) that TACC Stampede
    (Intel) gave 1e-8 core-hours per zone update, with near-perfect weak
    scaling up to 4096 cores.  Translation is 36e-6 core-seconds per zone.
    So, Torch with all its bells+whistles is about 4x slower.
    Not unexpected w/ radiative transfer, detailed heat/cool, wind deposition
    etc.

\item \textbf{Question:} when to do EOS update?  How?

    \textbf{Partial answer:}
    \url{http://flash.uchicago.edu/pipermail/flash-users/2017-May/002336.html}

\end{itemize}


% ============================
% Josh's old runs on Cartesius
% ============================
\section{Catalog of Josh's runs on Cartesius}

The Torch collaboration has project space on the supercomputer Cartesius,
located at the path \texttt{/lustre2/0/flashrad/}.

Josh has archived the following runs, here listed alphabetically with last-modified dates:
\begin{verbatim}
   4096 Jun 23  2018 DracoM3V02B							only "chk_0002"
 434176 Jun 28  2018 DracoM3V02BW
   4096 Aug 25  2018 M3V02AWL10_USM							empty
 147456 Aug 31  2018 M3V02AW_USM
 360448 Jan  1  2019 M3V02BWL10_USM
   4096 Aug 31  2018 M3V02BWL8_USM							empty

 344064 Aug  2  2018 M4V02A_cloud_evaporated
1261568 Mar 20  2019 M4V02AW
 540672 Feb 19  2019 M4V02AW_bad_IGIMF
 503808 Jul 24  2018 M4V02AWL10
   4096 Aug 25  2018 M4V02AWL10_USM
 786432 Jul 14  2018 M4V02AWL10_with_ndens_for_molcool_1e3
 544768 Aug 30  2018 M4V02AW_USM
 360448 Nov 21  2018 M4V02BWL10_USM
   4096 Aug 31  2018 M4V02BWL8_USM

 159744 Jun 23  2018 M5V02A
 270336 Oct 12  2018 M5V02AL8noFB_USM
   4096 Feb  7  2019 M5V02AW
 548864 Oct 18  2018 M5V02AWL10_USM
1175552 Feb 14  2019 M5V02AW_restart
 385024 Aug 27  2018 M5V02AW_USM
   4096 Sep  1  2018 M5V02BWL10_USM
   4096 Aug 31  2018 M5V02BWL8_USM

   4096 Mar 19  2019 M8e3_a08_x60_y370_z30					stratified box run from Juan Camilo Ibanez-Mejia

 196608 Jun 23  2018 oldDracoM3V02B
   4096 Mar  5  2019 radtest
      6 Jun 23  2018 test
   4096 Jun 23  2018 wrong_rad
\end{verbatim}

Several exist on the Drexel Univ cluster ``Draco'' as well.

Information about run wall-clock time used:

\begin{verbatim}
Here's an empirical measurement from one of Josh's 1e4 Msun runs (M4V02BWL10_USM).
It has 8 levels of refinement (lrefine_max=9), domain +/-24 pc, so max resolution 0.012 pc.
Note: this is NOT the run shown in Josh's paper.

    $ stat turbsph_hdf5_plt_cnt_2009
    $ h5dump -d "real scalars" turbsph_hdf5_plt_cnt_2009
        file modification time: 2018-11-21 04:50:04.000000000 +0100
        time = 1.06826e+14
        dt = 5.73224e+09

    $ stat turbsph_hdf5_plt_cnt_2011
    $ h5dump -d "real scalars" turbsph_hdf5_plt_cnt_2011
        file modification time: 2018-11-21 05:41:34.000000000 +0100
        time = 1.06887e+14
        dt = 6.04403e+09

This run used 8641 blocks of size 16^3, with 475 processors for hydro.
Advancing two plots (~2000 years) required ~50 wall-clock minutes and ~10 hydro steps.
One hydro step = ~5 wall-clock min, at 475 procs with ~18 blocks/core.
\end{verbatim}

Information about run disk usage:

\begin{verbatim}
Josh's 1e3-1e4 Msun runs reached, at the last plotted timestep in each run:

    Cloud   Plot    Plot    Plot        Chkpt   Chkpt   Chkpt     Last      Total   Run
    (Msun)  size    file#   interval    size    file#   interval  time      usage   directory
    ----------------------------------------------------------------------------------------------
    1e3     1.8G    2175    0.001 Myr   9.1G    71      0.1 Myr   5.4 Myr   1.7T    M3V02BWL10_USM  (Sep-Dec 2018)
    1e4     2.6G    2011    0.001 Myr   14G     360     0.01 Myr  3.4 Myr   4.4T    M4V02BWL10_USM  (Sep-Nov 2018)
    1e4     233M    1929    0.001 Myr   1.1G    3194    0.1 Myr   9.1 Myr   3.6T    M4V02AW         (Jun-Jul 2018)
\end{verbatim}

Notes:
disk space usage includes some analysis products/plots.
Plotfile sizes are adjustable (choose which FLASH fields to dump).
Particle file sizes are negligible (~1 MB or less).
Josh must have raised the plot frequency after the first star formed.

Tentatively, several to ~10 TB seems good for one production run at Josh's resolution, out to t=10 Myr.
On Cartesius, Josh applied for and got 400 TB.
As of October 2019, we are using 44 TB in total.

\end{document}
